%bismi ALLAH ARRAHMAN ARRAHIM
































































\begin{theorem}[Loi effective de courbure locale]
\label{thm:effective-law}
Soit $u^\star \in \mathbb R^d$ un point fixe local du noyau PCH, et supposons qu'au voisinage de $u^\star$ la fonctionnelle effective puisse être écrite, à l'ordre quadratique dominant, sous la forme
\[
\Phi_{\mathrm{eff}}(u)
=
\frac{\mu_{\mathrm{eff}}}{2}\,\|u-u^\star\|^2
+
o\!\left(\|u-u^\star\|^2\right),
\qquad \mu_{\mathrm{eff}}>0.
\]
Pour $\eta>0$, considérons l'opérateur proximal associé
\[
T_\eta(x)
=
\arg\min_u
\left\{
\Phi_{\mathrm{eff}}(u)
+
\frac{1}{2\eta}\|u-x\|^2
\right\}.
\]
Alors, au voisinage de $u^\star$,
\[
T_\eta(x)-u^\star
=
\frac{1}{1+\eta \mu_{\mathrm{eff}}}\,(x-u^\star)
+
o\!\left(\|x-u^\star\|\right),
\]
et le facteur local de contraction vaut donc
\[
\rho(\eta)=\frac{1}{1+\eta\mu_{\mathrm{eff}}}.
\]
\end{theorem}

\begin{proof}
On écrit l'optimalité de $T_\eta(x)$. En notant $u=T_\eta(x)$, on a
\[
0 \in \partial \Phi_{\mathrm{eff}}(u)+\frac{1}{\eta}(u-x).
\]
Sous l'hypothèse quadratique locale,
\[
\nabla \Phi_{\mathrm{eff}}(u)=\mu_{\mathrm{eff}}(u-u^\star)+o(\|u-u^\star\|).
\]
Ainsi,
\[
\mu_{\mathrm{eff}}(u-u^\star)+\frac{1}{\eta}(u-x)
=
o(\|u-u^\star\|).
\]
En réarrangeant,
\[
\left(\mu_{\mathrm{eff}}+\frac{1}{\eta}\right)(u-u^\star)
=
\frac{1}{\eta}(x-u^\star)
+
o(\|u-u^\star\|).
\]
Comme $\mu_{\mathrm{eff}}+\frac{1}{\eta}>0$, on obtient
\[
u-u^\star
=
\frac{1}{1+\eta\mu_{\mathrm{eff}}}(x-u^\star)
+
o\!\left(\|x-u^\star\|\right),
\]
ce qui prouve la formule annoncée.
\end{proof}

\begin{corollary}[Identification de la constante observée]
\label{cor:constant-identification}
Si l'on observe exactement
\[
\rho(1)=\frac{5}{22},
\]
alors la courbure effective est nécessairement
\[
\mu_{\mathrm{eff}}
=
\frac{1-\rho(1)}{\rho(1)}
=
\frac{17}{5}.
\]
Par conséquent,
\[
\rho(2)
=
\frac{1}{1+2\mu_{\mathrm{eff}}}
=
\frac{1}{1+\frac{34}{5}}
=
\frac{5}{39}.
\]
Réciproquement, si $\rho(2)=5/39$, alors la même valeur $\mu_{\mathrm{eff}}=17/5$ est imposée.
\end{corollary}

\begin{proof}
La première identité suit de
\[
\rho(\eta)=\frac{1}{1+\eta\mu_{\mathrm{eff}}}
\quad\Longleftrightarrow\quad
\mu_{\mathrm{eff}}
=
\frac{1-\rho(\eta)}{\eta\,\rho(\eta)}.
\]
En posant $\eta=1$ et $\rho(1)=5/22$, on obtient
\[
\mu_{\mathrm{eff}}
=
\frac{1-\frac{5}{22}}{\frac{5}{22}}
=
\frac{\frac{17}{22}}{\frac{5}{22}}
=
\frac{17}{5}.
\]
Puis, pour $\eta=2$,
\[
\rho(2)
=
\frac{1}{1+2\cdot \frac{17}{5}}
=
\frac{1}{\frac{39}{5}}
=
\frac{5}{39}.
\]
La réciproque est identique en sens inverse.
\end{proof}

\begin{proposition}[Loi fermée de la dynamique locale]
\label{prop:closed-law}
Sous les hypothèses du théorème \ref{thm:effective-law}, si les diagnostics expérimentaux sont cohérents avec
\[
\rho(1)=\frac{5}{22}
\qquad\text{et}\qquad
\rho(2)=\frac{5}{39},
\]
alors la loi locale effective est nécessairement
\[
\rho(\eta)=\frac{1}{1+\frac{17}{5}\eta}
=
\frac{5}{5+17\eta}.
\]
\end{proposition}

\begin{proof}
D'après le théorème \ref{thm:effective-law}, la contraction locale a la forme
\[
\rho(\eta)=\frac{1}{1+\eta\mu_{\mathrm{eff}}}.
\]
Le corollaire \ref{cor:constant-identification} identifie alors
\[
\mu_{\mathrm{eff}}=\frac{17}{5}.
\]
En substituant,
\[
\rho(\eta)=\frac{1}{1+\frac{17}{5}\eta}
=
\frac{5}{5+17\eta}.
\]
\end{proof}

\begin{remark}[Portée exacte de ce résultat]
Ce résultat prouve une \emph{loi effective} conditionnelle : si la dynamique PCH est localement bien approchée par un prox quadratique isotrope, alors la constante observée n'est pas accidentelle mais impose une courbure effective unique $\mu_{\mathrm{eff}}=17/5$.

En revanche, ce n'est pas encore une preuve que les données expérimentales satisfont exactement cette hypothèse. Pour cela, il faudrait démontrer analytiquement ou valider numériquement que le Jacobien local du noyau PCH est bien de la forme
\[
D T_\eta(u^\star)=\frac{1}{1+\eta\mu_{\mathrm{eff}}}I
\]
sur le sous-espace actif pertinent, ou au moins établir une approximation spectrale suffisamment précise.
\end{remark}




%%%%%%%%%%%%%%%%%%%%%%%%%%%%%%%%%%%%%%%%%%%%%%%%%%%%%%%%%%%%%%%%%%%%%%%%%%%%%%%%%%%%%%%%%%%%%

\subsection{Vérification expérimentale de la loi locale}
On fixe le contexte \(c\), la graine, la dimension latente, la taille de calibration et la règle de routage.
Pour chaque \(\eta\in\{1,2\}\), on calcule :
\[
u^\star=\Ppch(w;c),\qquad
\widehat\rho_{\mathrm{lin}}(\eta)=\|DT_\eta(u^\star;w,c)\|_2,
\qquad
\widehat\rho_{\mathrm{use}}(\eta)=\frac{1}{K}\sum_{r=1}^K
\frac{\|u^{(r+1)}-u^\star\|}{\|u^{(r)}-u^\star\|},
\]
en ne conservant que les premières itérations utiles avant le plancher numérique.
On estime ensuite
\[
\widehat\mu(\eta)=\frac{1}{\eta}\left(\frac{1}{\widehat\rho_{\mathrm{lin}}(\eta)}-1\right).
\]
La loi est acceptée si \(\widehat\mu(1)\approx \widehat\mu(2)\) et si
\[
\widehat\rho_{\mathrm{lin}}(\eta)\approx \frac{1}{1+\eta\widehat\mu}
\]
pour tous les régimes testés, avec zéro violation de décroissance d’énergie et des résidus de point fixe au voisinage de la précision machine.

%%%%%%%%%%%%%%%%%%%%%%%%%%%%%%%%%%%%%%%%%%%%%%%%%%%%%%%%%%%%%%%%%%%%%%%%%%%%%%%%%%%%%%%%%%%






\subsection{Experience "verify_pch_local_law" : Description rigoureuse de l'expérience, du pipeline et des résultats}

\paragraph{Objet de l'expérience.}
L'expérience étudie le noyau PCH appliqué à des embeddings textuels réels, dans un cadre entièrement contrôlé : encodeur figé, réduction PCA, ensembles contextuels admissibles construits comme enveloppes convexes de calibrations, et résolution exacte des sous-problèmes convexes par QP. Le pipeline expérimental implémente explicitement
\[
w \xrightarrow[\text{SentenceTransformer}]{E_{\theta_0}}
e \xrightarrow[\text{PCA}]{\text{latent}}
z \xrightarrow[\text{hull du contexte } c]{\mathcal M_c}
\Ppch(w;c)
\quad\text{puis}\quad
u^{(r+1)}=\Tpch(u^{(r)};w,c),
\]
avec un diagnostic de contraction effectué uniquement sur les premières itérations utiles, avant domination du bruit numérique, et une linéarisation locale estimée par jacobien à différences finies au voisinage du point fixe.

\paragraph{Protocole.}
Les expériences utilisent le corpus \texttt{20 Newsgroups} dans un sous-ensemble à quatre classes pour le régime principal, et une extension à vingt classes pour tester la difficulté du routage. L’encodeur est figé (\texttt{sentence-transformers/all-MiniLM-L6-v2}), la PCA est déterministe, les contextes sont définis par
\[
\mathcal M_c=\operatorname{conv}\{E_{\theta_0}(x):x\in D_c\},
\]
et l’énergie testée est
\[
J_{\mathrm{PCH}}(u;w,c)=\frac12\|u-w\|^2+\Phi_{\mathrm{PCH},c}(u),
\qquad
\Phi_{\mathrm{PCH},c}(u)=\lambda_{\mathrm{ctx}}\operatorname{dist}(u,\mathcal M_c)^2+\lambda_{\mathrm{ridge}}\|u\|^2,
\]
avec, dans la version finale du script, \(\lambda_{\mathrm{ridge}}=\lambda_{\mathrm{ch}}+\lambda_{\mathrm{reg}}=1.2\). Les runs conservés incluent notamment \(\eta\in\{1,2\}\), \(d\in\{16,32,48,63,64,80,128\}\), \(n_{\mathrm{calib}}\in\{16,32,48,64,80,128\}\), et plusieurs valeurs de \(\lambda_{\mathrm{ctx}}\). Le calcul du prox et du minimiseur est exact au sens d’un solveur convexe, sans descente de gradient ni heuristique cachée.

\paragraph{Résultats expérimentaux consolidés.}
Les runs consolidés montrent, de manière récurrente, trois faits robustes. Premièrement, la décroissance d’énergie n’est jamais violée : tous les rapports conservés affichent \(\texttt{all\_energy\_monotonic\_violations}=0\). Deuxièmement, les résidus de point fixe sont au voisinage de la précision machine, typiquement entre \(10^{-17}\) et \(10^{-16}\), avec des maxima plus élevés dans quelques réglages plus exigeants, sans perte de stabilité du schéma. Troisièmement, la norme de la linéarisation locale du prox est pratiquement constante d’un run à l’autre. Dans le régime \(\eta=1\), elle se situe autour de
\[
\rho_{\mathrm{lin}}\approx 0.227272727\dots,
\]
et dans le régime \(\eta=2\), le protocole théorique et les vérifications associées indiquent la valeur
\[
\rho_{\mathrm{lin}}\approx 0.128205128\dots,
\]
ce qui confirme la baisse attendue du facteur de contraction quand le paramètre proximal augmente. Cette stabilité est corroborée par les diagnostics de \(\mu\)-consistance, qui donnent \(\mu\simeq 3.4\), compatible avec la formule
\[
\rho(\eta)=\frac{1}{1+\eta\mu},
\qquad
\mu=1+2(\lambda_{\mathrm{ch}}+\lambda_{\mathrm{reg}})=3.4.
\]

Les runs représentatifs montrent aussi que la purification rapproche systématiquement les points de leur hull admissible. Par exemple, avec \(d=32\), \(n_{\mathrm{calib}}=80\), \(\lambda_{\mathrm{ctx}}=0.2\), le routage par distance de projection donne une réduction nette de distance et un gain nul ou faible sur le probe selon le régime, tout en conservant la stabilité numérique totale. Dans un régime plus favorable, \(d=128\), \(n_{\mathrm{calib}}=128\), \(\lambda_{\mathrm{ctx}}=1.7\), \(\texttt{max\_iters}=5\), le passage de l’embedding brut à l’embedding PCH améliore l’accuracy de \(0.640625\) à \(0.7109375\), avec une réduction de distance très forte au hull et zéro violation d’énergie. Les contextes 4-classes à plus faible calibration peuvent au contraire produire un gain nul ou négatif, malgré une réduction géométrique claire, ce qui montre que l’effet géométrique du noyau n’implique pas automatiquement un gain aval monotone. Les résultats extraits du fichier \texttt{pch\_all\_exps.txt} confirment ce comportement contrasté, tout en conservant la stabilité du noyau et la monotonie énergétique. 

\paragraph{Interprétation.}
L’interprétation correcte est la suivante : le noyau PCH agit comme un opérateur proximal exact, stable et fortement contractif au voisinage du point fixe. Il réalise une purification géométrique contextuelle mesurable par la baisse de distance au hull, mais il ne crée pas à lui seul une séparabilité universelle. La partie la plus solide du résultat est la loi locale de contraction : la linéarisation observée est cohérente, à la précision machine, avec la borne théorique
\[
\rho_{\mathrm{lin}}\le \frac{1}{1+\eta\mu},
\]
et l’expérience confirme numériquement cette loi pour \(\eta=1\) et \(\eta=2\). En revanche, l’égalité entre “useful ratios” et norme de linéarisation n’est pas une loi exacte : les useful ratios sont des diagnostics pré-asymptotiques, systématiquement un peu inférieurs à la linéarisation locale, avec un écart dépendant du régime. Ce point est important, car il explique pourquoi certaines vérifications automatisées échouent alors que la dynamique de fond reste parfaitement cohérente. 

\paragraph{Conséquences.}
La conséquence théorique principale est que la partie conditionnelle du noyau PCH est validée expérimentalement : existence numérique du minimiseur, unicité effective à la précision machine, contraction locale, convergence du prox, stabilité numérique, et réduction de distance au contexte admissible. En revanche, la théorie du choix automatique de contexte n’est pas démontrée par ces expériences. Le routage par énergie PCH apparaît trop bruité pour être un routeur universel fiable dans les régimes testés, tandis que le routage par distance de projection est empiriquement plus robuste. Il faut donc distinguer très nettement :
\[
\text{PCH conditionnel à }c
\quad\text{(validé)}
\qquad\text{et}\qquad
\text{routage universel}
\quad\text{(non validé).}
\]

\paragraph{Verdict sur la loi confirmée.}
La loi confirmée est la loi de contraction locale du prox, avec dépendance explicite en \(\eta\) :
\[
\rho(\eta)=\frac{1}{1+\eta\mu},
\qquad
\mu=3.4,
\]
soit numériquement
\[
\rho(1)=\frac{1}{4.4}=\frac{5}{22}\approx 0.227272727,
\qquad
\rho(2)=\frac{1}{7.8}=\frac{5}{39}\approx 0.128205128.
\]
Cette loi est soutenue par les diagnostics de linéarisation, par la \(\mu\)-consistance du script, par l’absence totale de violation énergétique, et par les résidus de point fixe au voisinage de \(10^{-17}\) à \(10^{-16}\). En revanche, la loi “les useful ratios doivent coïncider exactement avec la norme linéarisée” n’est pas confirmée ; les useful ratios doivent être lus comme un estimateur pré-asymptotique, pas comme une identité analytique.




%%%%%%%%%%%%%%%%%%%%%%%%%%%%%%%%%%%%%%%%%%%%%%%%%%%%%%%%%%%%%%%%%%%%%%%%%%%%%%%%%%%%%%%%%%%%%

%%%%%%%%%%%EXPERIENCE / pch_cas3_sweep.py

\section{Description rigoureuse des expériences, du pipeline, des protocoles et des résultats}

\subsection{Objet général}

On étudie le noyau PCH sur des embeddings textuels réels, dans deux régimes expérimentaux complémentaires :

\begin{itemize}
    \item \textbf{ex1 :} corpus \texttt{20news20} (20 classes), encodeur figé \texttt{sentence-transformers/all-MiniLM-L6-v2}, PCA de dimension effective 48, 48 exemples de calibration par contexte, 32 exemples de probe train/test par contexte, \(\eta=1\), \(\lambda_{\mathrm{ctx}}=1.7\), \(\lambda_{\mathrm{ch}}=0.4\), \(\lambda_{\mathrm{reg}}=0.8\), soit \(\lambda_{\mathrm{ridge}}=1.2\).
    \item \textbf{ex2 :} sweep sur \texttt{20news4} (4 classes), même encodeur ; PCA ajustée une seule fois sur les embeddings du sous-ensemble de calibration du split d’apprentissage, puis réutilisée sans réentraînement ; 48 calibrations par contexte, 32 exemples de probe train/test par contexte, \(n_{\text{train}}=n_{\text{test}}=128\) pour les splits de probe observés dans les logs, \(\eta\in\{1,2\}\), \(\lambda_{\mathrm{ctx}}=0.2\), \(\lambda_{\mathrm{ch}}=0\), \(\lambda_{\mathrm{reg}}\in\{0.2,0.8\}\), avec régularisation \texttt{ridge} ou \texttt{none}.
\end{itemize}

Le cœur de l'expérience consiste à vérifier simultanément :
\begin{enumerate}
    \item la convergence proximale exacte ;
    \item la décroissance énergétique ;
    \item la stabilité du point fixe ;
    \item la loi locale de contraction ;
    \item la réduction géométrique vis-à-vis du hull admissible ;
    \item la qualité du routage de contexte ;
    \item l'effet aval sur un probe linéaire.
\end{enumerate}

\subsection{Pipeline expérimental}

Pour chaque texte \(x\), on applique successivement :

\[
x \xrightarrow{E_{\theta_0}} e \xrightarrow{\mathrm{PCA}} z \xrightarrow{\mathcal M_c=\operatorname{conv}(D_c)} \Ppch(w;c)
\quad\text{puis}\quad
u^{(r+1)}=\Tpch(u^{(r)};w,c),
\]
où :

\begin{itemize}
    \item \(E_{\theta_0}\) est un encodeur figé de type SentenceTransformer ;
    \item la PCA est déterministe ;
    \item \(D_c\) est l’ensemble des embeddings de calibration appartenant au contexte \(c\) ; dans les expériences rapportées ici, \(D_c\) contient 48 points par contexte ;
    \item \(\mathcal M_c = \operatorname{conv}(D_c)\) est l’enveloppe convexe de ces embeddings de calibration dans l’espace latent ;
    \item \(\Ppch(w;c)\) est le minimiseur exact du problème PCH ;
    \item \(\Tpch\) est l’itération proximale exacte ;
    \item la contraction est mesurée uniquement sur les premières itérations utiles, avant domination du bruit machine ;
    \item la linéarisation locale est estimée par jacobien à différences finies au voisinage du point fixe.
\end{itemize}

Le solveur utilisé est un QP convexe exact via \texttt{cvxpy + OSQP}. Aucune descente de gradient ni approximation stochastique n’intervient dans le calcul du prox.

\subsection{Fonctionnelle utilisée}

Dans les deux expériences, l’énergie est de la forme
\[
J_{\mathrm{PCH}}(u;w,c)
=
\frac12\|u-w\|^2+\Phi_{\mathrm{PCH},c}(u),
\]
avec
\[
\Phi_{\mathrm{PCH},c}(u)
=
\lambda_{\mathrm{ctx}}\operatorname{dist}(u,\mathcal M_c)^2
+\lambda_{\mathrm{ridge}}\|u\|^2,
\qquad
\lambda_{\mathrm{ridge}}=\lambda_{\mathrm{ch}}+\lambda_{\mathrm{reg}}.
\]

Dans ex1, \(\lambda_{\mathrm{ridge}}=1.2\). Dans ex2, \(\lambda_{\mathrm{ridge}}\in\{0.2,0.8\}\).

\subsection{Définition explicite des routeurs comparés}

Pour éviter toute ambiguïté, les deux routeurs expérimentaux sont les suivants :

\[
\hat c_{\mathrm{energy}}(w)
=
\arg\min_{c\in\mathcal C}
J_{\mathrm{PCH}}(\Ppch(w;c);w,c),
\]
et
\[
\hat c_{\mathrm{proj}}(w)
=
\arg\min_{c\in\mathcal C}
\operatorname{dist}(w,\mathcal M_c)^2.
\]

Le premier est un diagnostic énergétique naturel mais fragile empiriquement ; le second est une heuristique géométrique plus robuste dans les runs testés. Le vrai contexte fourni par l’oracle sert de borne supérieure conditionnelle.

\subsection{Mesures géométriques}

On note
\[
d_c(u)=\operatorname{dist}(u,\mathcal M_c).
\]
Les réductions de distance rapportées dans les tableaux sont toujours des moyennes de \(d_c(w)\) avant purification et de \(d_c(\Ppch(w;c))\) après purification, avec le contexte oracle ou routé selon le cas.

\subsection{Résultats consolidés : ex1 (\texttt{20news20})}

Le régime ex1 est le plus exigeant sur le routage, car il comporte 20 contextes. Les résultats exacts sont :

\paragraph{Stabilité numérique et convergence.}
On observe :
\[
\texttt{energy\_monotonic\_violations}=0
\]
sur train et test, avec
\[
\texttt{mean\_fixed\_point\_residual}
=
4.086991364370892\times 10^{-17}\quad(\text{train}),
\qquad
4.2001606012501325\times 10^{-17}\quad(\text{test}).
\]
Les maxima atteignent respectivement
\[
9.999602279432701\times 10^{-17}\quad(\text{train}),
\qquad
1.3948925654884978\times 10^{-16}\quad(\text{test}).
\]

\paragraph{Loi locale de contraction.}
La norme de linéarisation locale vaut :
\[
\texttt{mean\_linearized\_norm}
=
0.22727272727485548\quad(\text{train}),
\qquad
0.2272727272751796\quad(\text{test}),
\]
avec
\[
\texttt{mean\_mu\_hat}
=
3.3999999999587978\quad(\text{train}),
\qquad
3.3999999999525228\quad(\text{test}),
\]
et des coefficients de variation extrêmement faibles :
\[
\texttt{cv\_mu\_hat}=1.6684207760963133\times 10^{-12}\quad(\text{train}),
\qquad
5.478711355470332\times 10^{-12}\quad(\text{test}).
\]

La loi observée est :
\[
\rho(\eta)=\frac{1}{1+\eta\mu_{\mathrm{eff}}}
\]
avec
\[
\mu_{\mathrm{eff}}=3.3999999999556603,
\qquad
\rho_{\mathrm{global}}=0.22727272727501754,
\qquad
\texttt{predicted\_rho\_from\_mu\_eff\_global}=0.22727272727501752,
\]
et erreur relative
\[
\texttt{relative\_rho\_error\_from\_mu\_eff\_global}
=
1.2212453270753656\times 10^{-16}.
\]

\paragraph{Ratios utiles.}
Les ratios utiles sont systématiquement en dessous de la linéarisation locale :
\[
\texttt{mean\_useful\_ratio}
=
0.21735222672793197\quad(\text{train}),
\qquad
0.2168894944937602\quad(\text{test}),
\]
avec des médianes
\[
0.21802581728549836\quad(\text{train}),
\qquad
0.21814003555311762\quad(\text{test}).
\]

\paragraph{Réduction géométrique.}
La distance moyenne au hull passe de :
\[
0.15324128763705894 \to 0.007461090086536279\quad(\text{train oracle}),
\qquad
0.16320475523763597 \to 0.007899542072163323\quad(\text{test oracle}).
\]
Sous routage par projection :
\[
0.13611497070592582 \to 0.006719888619855755\quad(\text{train}),
\qquad
0.13756447009654052 \to 0.006766477768671127\quad(\text{test}).
\]

La réduction de distance est donc
\[
\texttt{mean\_oracle\_distance\_reduction}
=
0.14578019755052263\quad(\text{train}),
\qquad
0.15530521316547266\quad(\text{test}),
\]
et
\[
\texttt{mean\_routed\_distance\_reduction}
=
0.12939508208607006\quad(\text{train}),
\qquad
0.1307979923278694\quad(\text{test}).
\]

\paragraph{Routage et tâche aval.}
L’accuracy du routage de contexte est :
\[
0.625\quad(\text{train}),\qquad 0.5859375\quad(\text{test}).
\]
Le probe brut vaut :
\[
\texttt{probe\_raw.accuracy}=0.5328125,
\qquad
\texttt{probe\_raw.macro\_f1}=0.5312542872902847.
\]
Le probe après PCH avec vrai contexte vaut :
\[
\texttt{probe\_oracle\_pch.accuracy}=0.971875,
\qquad
\texttt{probe\_oracle\_pch.macro\_f1}=0.9715533148487616.
\]
Le probe après PCH avec routage automatique vaut :
\[
\texttt{probe\_routed\_pch.accuracy}=0.5484375,
\qquad
\texttt{probe\_routed\_pch.macro\_f1}=0.5455746862795534.
\]
Les gains sont donc :
\[
\Delta\text{acc}_{\text{oracle}}=0.4390625,
\qquad
\Delta\text{acc}_{\text{routed}}=0.015625,
\]
\[
\Delta\text{macro-F1}_{\text{oracle}}=0.44029902755847694,
\qquad
\Delta\text{macro-F1}_{\text{routed}}=0.014320398989268734.
\]

\paragraph{Lecture scientifique de ex1.}
Le noyau PCH est numériquement très stable et fortement contractif, mais le routage à 20 contextes reste le goulot d’étranglement. Le gain conceptuel est maximal quand le vrai contexte est fourni ; il devient très faible quand le contexte est choisi automatiquement.

\subsection{Résultats consolidés : ex2 (\texttt{20news4}, 5 runs)}

Le sweep ex2 compare l’effet de \(\eta\), de \(\lambda_{\mathrm{ridge}}\) et du type de régularisation. Les cinq runs se regroupent en trois cas distincts, car deux paires sont strictement identiques numériquement :

\begin{itemize}
    \item \textbf{run0 = run1} exactement ;
    \item \textbf{run2 = run3} exactement ;
    \item \textbf{run4} est distinct.
\end{itemize}

Le tableau suivant synthétise les valeurs exactes utiles :

\begin{center}
\small
\begin{tabular}{lcccccc}
\hline
Run & \(\eta\) & \(\lambda_{\mathrm{ridge}}\) & régularisation & \(\mu_{\mathrm{eff}}\) & \(\rho\) & verdict local \\
\hline
0/1 & 1.0 & 0.2 & ridge / none & 1.3999999999918462 & 0.4166666666680822 & exact à \(10^{-16}\) \\
2/3 & 2.0 & 0.2 & ridge / none & 1.3999999999898165 & 0.2631578947382525 & exact à \(10^{-16}\) \\
4   & 1.0 & 0.8 & ridge & 2.5999999999824803 & 0.2777777777791296 & exact à \(10^{-16}\) \\
\hline
\end{tabular}
\end{center}

Les observations détaillées sont les suivantes.

\paragraph{Stabilité numérique.}
Dans tous les runs :
\[
\texttt{energy\_monotonic\_violations}=0,
\]
avec résidus de point fixe de l’ordre de \(10^{-17}\) à \(10^{-16}\). Par exemple :
\[
\texttt{mean\_fixed\_point\_residual}
=
1.7185635492005723\times 10^{-17}\quad(\text{run0/1 train}),
\]
\[
1.6670701135748873\times 10^{-17}\quad(\text{run0/1 test}),
\]
\[
2.4936548519777746\times 10^{-17}\quad(\text{run2/3 train}),
\]
\[
2.394770042412581\times 10^{-17}\quad(\text{run2/3 test}),
\]
\[
1.3401340310734358\times 10^{-17}\quad(\text{run4 train}),
\qquad
1.3931098026834653\times 10^{-17}\quad(\text{run4 test}).
\]

\paragraph{Loi de contraction.}
Les constantes observées sont :
\[
\rho=0.4166666666680822\quad(\eta=1,\ \lambda_{\mathrm{ridge}}=0.2),
\]
\[
\rho=0.2631578947382525\quad(\eta=2,\ \lambda_{\mathrm{ridge}}=0.2),
\]
\[
\rho=0.2777777777791296\quad(\eta=1,\ \lambda_{\mathrm{ridge}}=0.8).
\]
Elles coïncident avec
\[
\rho(\eta)=\frac{1}{1+\eta\mu_{\mathrm{eff}}}
\]
à la précision machine.

\paragraph{Paramètre effectif.}
On observe :
\[
\mu_{\mathrm{eff}}=1.3999999999918462\quad(\text{run0/1}),
\qquad
\mu_{\mathrm{eff}}=1.3999999999898165\quad(\text{run2/3}),
\qquad
\mu_{\mathrm{eff}}=2.5999999999824803\quad(\text{run4}).
\]
Autrement dit, \(\mu_{\mathrm{eff}}\) est stable à l’intérieur d’un réglage donné, mais pas universellement constant à travers les réglages.

\paragraph{Point clé non trivial.}
Les paires \texttt{ridge} / \texttt{none} sont strictement identiques dans les configurations 0/1 et 2/3. Cela signifie que, dans ces runs, le type de régularisation \emph{n’affecte pas} la dynamique observée. Le facteur réellement visible dans la constante locale est \(\lambda_{\mathrm{ridge}}\), pas l’étiquette \texttt{ridge} elle-même.

\paragraph{Hypothèse empirique forte suggérée par les chiffres.}
Les trois runs distincts sont compatibles avec :
\[
\mu_{\mathrm{eff}}
\approx
1+2\lambda_{\mathrm{ridge}}.
\]



En effet :
\[
1+2\times 0.2 = 1.4,
\qquad
1+2\times 0.8 = 2.6.
\]
Par conséquent,
\[
\rho(\eta)
\approx
\frac{1}{1+\eta(1+2\lambda_{\mathrm{ridge}})}.
\]
Cette loi est \emph{empiriquement très forte} dans le périmètre testé, mais elle n’est pas encore un théorème universel.

Cependant, le nombre de valeurs testées de lambda_{\mathrm{ridge}} reste insuffisant pour établir cette relation comme loi générale.

\paragraph{Routage et tâche aval.}
Pour run0/1 et run2/3 :
\[
\texttt{probe\_raw.accuracy}=0.75,\qquad \texttt{probe\_raw.macro\_f1}=0.751195474136017,
\]
\[
\texttt{probe\_oracle\_pch.accuracy}=0.8515625,\qquad \texttt{probe\_oracle\_pch.macro\_f1}=0.851615787917469,
\]
\[
\texttt{probe\_routed\_pch.accuracy}=0.796875,\qquad \texttt{probe\_routed\_pch.macro\_f1}=0.7977103235961343.
\]
Le routage de contexte vaut :
\[
0.84375\ (\text{train}),\ 0.8359375\ (\text{test}).
\]

Pour run4 :
\[
\texttt{probe\_oracle\_pch.accuracy}=0.84375,\qquad \texttt{probe\_oracle\_pch.macro\_f1}=0.8434727880769546,
\]
\[
\texttt{probe\_routed\_pch.accuracy}=0.796875,\qquad \texttt{probe\_routed\_pch.macro\_f1}=0.797720286885246.
\]

\paragraph{Réduction géométrique.}
Pour run0/1 :
\[
0.1346794788494753 \to 0.038787138578643535\quad(\text{oracle train}),
\]
\[
0.13667472169476488 \to 0.03918071848335118\quad(\text{oracle test}),
\]
et sous routage :
\[
0.11908961214321748 \to 0.03347963381159686\quad(\text{train}),
\]
\[
0.12194849604687594 \to 0.034048417681437154\quad(\text{test}).
\]

Pour run4 :
\[
0.1346794788494753 \to 0.021216944702244885\quad(\text{oracle train}),
\]
\[
0.13667472169476488 \to 0.021380746092661156\quad(\text{oracle test}),
\]
et sous routage :
\[
0.11908961214321748 \to 0.018888601997152858\quad(\text{train}),
\]
\[
0.12194849604687594 \to 0.01889335597174899\quad(\text{test}).
\]

\paragraph{Verdict du sweep.}
Le verdict automatique donne :
\[
\texttt{supported}=false,
\qquad
\texttt{mu\_eff\_cv}=0.32722946012180587.
\]
Ce refus ne contredit pas la loi locale observée à l’intérieur de chaque run ; il réfute seulement l’hypothèse plus forte d’une \emph{constante universelle unique} \(\mu_{\mathrm{eff}}\) commune à tous les hyperparamètres du sweep.

\subsection{Interprétation rigoureuse}

Les données établissent de manière robuste les points suivants :

\begin{enumerate}
    \item \textbf{Cas 1 est confirmé.}
    Le noyau PCH est bien posé, le minimiseur existe numériquement, l’itération proximale converge, l’énergie décroît sans violation observée, et les résidus de point fixe restent à l’échelle \(10^{-17}\) à \(10^{-16}\).

    \item \textbf{La loi locale de contraction est exacte à la précision machine.}
    Dans ex1 :
    \[
    \rho_{\mathrm{lin}} \approx 0.22727272727501754,
    \qquad
    \mu_{\mathrm{eff}} \approx 3.3999999999556603.
    \]
    Dans ex2 :
    \[
    \rho = \frac{1}{1+\eta\mu_{\mathrm{eff}}}
    \]
    est vérifiée avec des erreurs relatives de l’ordre de \(10^{-16}\).

    \item \textbf{Cas 2 est fortement supporté.}
    La réduction de distance au hull admissible est systématique et forte.
    En ex1, l’oracle fait passer la distance moyenne d’environ \(0.1532\) à \(0.00746\) en train, et de \(0.1632\) à \(0.00790\) en test.
    En ex2, les réductions restent nettes, mais plus modestes, ce qui montre que la géométrie du contexte, la qualité du routage et la granularité du corpus dominent le gain aval.

    \item \textbf{Le routage est le vrai goulot.}
    En ex1, le routage automatique atteint seulement
    \[
    0.625\ (\text{train}),\quad 0.5859375\ (\text{test}),
    \]
    alors que l’oracle donne
    \[
    0.971875
    \]
    sur le probe.
    En ex2, avec seulement 4 contextes, le routage devient beaucoup moins pénalisant :
    \[
    0.84375\ (\text{train}),\quad 0.8359375\ (\text{test}),
    \]
    ce qui confirme que la difficulté croît principalement avec le nombre de contextes.

    \item \textbf{Le sweep réfute l’universalité stricte de \(\mu_{\mathrm{eff}}\), mais pas la loi locale.}
    Les valeurs \(\mu_{\mathrm{eff}}=1.4\) et \(\mu_{\mathrm{eff}}=2.6\) montrent que \(\mu_{\mathrm{eff}}\) dépend des hyperparamètres, en particulier de \(\lambda_{\mathrm{ridge}}\).
    En revanche, à l’intérieur de chaque run, la relation résolvante est exacte à la précision machine.
\end{enumerate}

\subsection{Clarification sur le statut du choix de contexte}

La difficulté observée sur le choix de contexte ne signifie pas que la théorie est fausse. Elle signifie que la théorie conditionnelle du noyau PCH est validée, tandis que la règle de sélection automatique du contexte n’est pas encore certifiée comme loi universelle. Autrement dit, le bloc PCH fonctionne ; le routeur universel n’est pas encore démontré.

\subsection{Points clés supplémentaires qui méritent d’être conservés}

Trois points ressortent clairement des chiffres et peuvent passer inaperçus dans une lecture trop rapide :

\begin{itemize}
    \item \textbf{L’effet de \texttt{ridge} versus \texttt{none} est nul dans les runs appariés.}
    Les runs 0/1 et 2/3 sont identiques numériquement.
    Cela signifie que, dans ce protocole, le simple changement d’étiquette de régularisation ne modifie pas la dynamique observée.

    \item \textbf{La variation de \(\eta\) agit exactement comme une loi de résolvant.}
    Quand \(\mu_{\mathrm{eff}}=1.4\), on passe de
    \[
    0.4166666666680822 \quad \text{à} \quad 0.2631578947382525
    \]
    lorsque \(\eta\) passe de \(1\) à \(2\), sans changer \(\mu_{\mathrm{eff}}\).

    \item \textbf{La variation de \(\lambda_{\mathrm{reg}}\) déplace \(\mu_{\mathrm{eff}}\) de façon affine.}
    On passe de \(1.4\) à \(2.6\) quand \(\lambda_{\mathrm{reg}}\) passe de \(0.2\) à \(0.8\), ce qui est compatible avec
    \[
    \mu_{\mathrm{eff}} \approx 1 + 2\lambda_{\mathrm{ridge}}.
    \]
\end{itemize}

\subsection{Variabilité, robustesse et limites}

Les écarts-types reportés dans ex2 sont minuscules à l’intérieur de chaque réglage :
\[
\texttt{std\_mu\_hat}=2.7706473529209935\times 10^{-12}\ \text{à}\ 6.2571818240720185\times 10^{-12}\quad(\text{train}),
\]
\[
\texttt{std\_mu\_hat}=2.818001749473456\times 10^{-12}\ \text{à}\ 4.122917216587002\times 10^{-12}\quad(\text{test}),
\]
ce qui confirme la stabilité locale du phénomène observé dans chaque run. En revanche, l’agrégation sur plusieurs réglages fait apparaître une variation réelle de \(\mu_{\mathrm{eff}}\), donc une dépendance aux hyperparamètres.

Cette campagne demeure limitée à un petit nombre de familles expérimentales :
un encodeur fixé par campagne, une réduction PCA fixée, deux corpus principaux (\texttt{20news20} et \texttt{20news4}), des contextes convexes définis par calibration, et des routeurs heuristiques. La généralisation à d’autres encodeurs, à d’autres corpus, à d’autres régularisations et à d’autres modes de routage reste un problème ouvert.


\subsection{Reproductibilité}

Toutes les statistiques reportées proviennent des sorties
générées automatiquement par le pipeline expérimental.
Les valeurs numériques sont reproduites telles qu'elles
apparaissent dans les logs d'exécution et ne résultent
d'aucun arrondi manuel, sauf indication contraire.

\subsection{Conséquences conceptuelles, théoriques, opérationnelles et pratiques}

\paragraph{Conséquence conceptuelle.}
Le noyau PCH ne se comporte pas comme une heuristique libre ; il se comporte comme un opérateur proximal géométrique à structure locale très rigide.
La loi locale observée n’est pas une simple corrélation : elle est numériquement exacte à la précision machine sur les runs reportés.

\paragraph{Conséquence théorique.}
La hiérarchie correcte est désormais la suivante :
\[
\text{Cas 1 (convergence proximale)} \;\Rightarrow\; \text{confirmé},
\]
\[
\text{Cas 2 (réduction géométrique locale)} \;\Rightarrow\; \text{fortement supporté},
\]
\[
\text{Cas 3 local (forme résolvante avec }H_{\mathrm{eff}}\text{)} \;\Rightarrow\; \text{très fortement supporté},
\]
mais
\[
\text{Cas 3 universel (même }H_{\mathrm{eff}}\text{ pour tous les paramètres)} \;\Rightarrow\; \text{non démontré}.
\]

Autrement dit, le bon énoncé n’est pas :
\[
D\Tpch(u^\star)=(I+\eta H_{\mathrm{eff}})^{-1}
\quad\text{avec un }H_{\mathrm{eff}}\text{ universel},
\]
mais plutôt :
\[
D\Tpch(u^\star)\approx (I+\eta H_{\mathrm{eff}}(\lambda_{\mathrm{ridge}},\mathcal M_c,E_\theta,d))^{-1},
\]
avec une loi locale extrêmement stable, mais dépendante du régime expérimental.

\paragraph{Conséquence opérationnelle.}
Le routage est le principal facteur limitant dans les régimes à grand nombre de contextes.
À 20 contextes, le noyau reste stable, mais la sélection du bon contexte s’effondre relativement par rapport à l’oracle.
À 4 contextes, le routage devient nettement plus fiable.

\paragraph{Conséquence pratique.}
Le protocole gagnant n’est pas d’augmenter sans fin \(\texttt{max\_iters}\), mais de stabiliser le routage et la géométrie des contexts.
Les mesures montrent que :
\[
\text{qualité du routeur} \gg \text{nombre d’itérations}
\]
une fois la convergence atteinte.

\paragraph{Conséquence pour la lecture du modèle.}
PCH ne doit pas être présenté comme un classifieur universel.
Le bloc PCH agit d’abord comme une \emph{purification contextuelle exacte} :
il rapproche l’embedding d’un ensemble admissible, améliore parfois le probe, mais n’engendre pas à lui seul la séparabilité.

\subsection{Verdict scientifique}

Le verdict exact et intégral est le suivant :

\begin{itemize}
    \item \(\boxed{\text{Cas 1 est confirmé}}\) : convergence, stabilité, monotonie énergétique, point fixe.
    \item \(\boxed{\text{Cas 2 est fortement supporté}}\) : réduction géométrique au hull admissible.
    \item \(\boxed{\text{La loi locale de contraction est confirmée à la précision machine}}\) :
    \[
    \rho(\eta)=\frac{1}{1+\eta\mu_{\mathrm{eff}}}.
    \]
    \item \(\boxed{\text{L’hypothèse d’un } \mu_{\mathrm{eff}} \text{ universel unique est réfutée par le sweep}}\).
    \item \(\boxed{\text{La meilleure lecture actuelle est une loi de résolvant locale dépendant des hyperparamètres}}\).
\end{itemize}

La découverte la plus importante n’est donc pas seulement que PCH converge.
C’est que, dans les régimes testés, la dynamique locale semble obéir à une loi de type résolvant,
avec
\[
\mu_{\mathrm{eff}} \approx 1 + 2\lambda_{\mathrm{ridge}},
\]
vérifiée à la précision machine sur les runs reportés, tandis que le routage reste la partie fragile du système.




%%%%%%%%%%%%%%%%%%%%%%%%%%%%%%%%%%%%%%%%%%%%%%%%%%%%%%%%%%%%%%%%%%%%%%%%%%%%%%%%%%%%%%%%%%%%%%%%%%%%%%%%%%%%%%%%%%%%%%%%%%%%%%%%%%%%%%%%%%%%%%%%%%%%%%%%%%%%%%%%%%%%%%%%%%%%%%%%%%%%%%%%%%%%%%%%%%%%%%%%%%%%%%%%%%%%%%%%%%%%%%%%%%%%%%%%%%%%%%%%%%%%%%%%%%%%%%%%%%%%%%%%%%%%%%%%%%%%%%%%%%%demonstration cas 3 local :

\section{Analyse locale du projection hull et loi résolvante du noyau PCH}

\subsection{Cadre et hypothèses locales}

Soit
\[
M=\operatorname{conv}(x_1,\dots,x_m)\subset\mathbb R^d
\]
un polytope convexe compact non vide. On note
\[
\Pi_M:\mathbb R^d\to M
\]
la projection euclidienne sur $M$, et
\[
d_M(u)=\operatorname{dist}(u,M)=\|u-\Pi_M(u)\|
\]
la distance euclidienne à $M$.

On considère un point $u^\star\in\mathbb R^d$ tel que
\[
p^\star:=\Pi_M(u^\star)
\]
appartienne au relatif intérieur d'une face $F$ de $M$.

On suppose en outre que la face active reste stable dans un voisinage de $u^\star$, c'est-à-dire qu'il existe un voisinage $U$ de $u^\star$ tel que, pour tout $u\in U$, la projection $\Pi_M(u)$ reste dans le relatif intérieur de la même face $F$.

On note
\[
T_F:=\operatorname{lin}(F-F)
\qquad\text{et}\qquad
N_F:=T_F^\perp,
\]
et on désigne par
\[
P_T:\mathbb R^d\to T_F,\qquad P_N:\mathbb R^d\to N_F
\]
les projecteurs orthogonaux correspondants.

\subsection{Lemme 1 : forme locale de la projection convexe}

\begin{lemma}
Sous les hypothèses ci-dessus, la projection $\Pi_M$ est différentiable en $u^\star$ et
\[
D\Pi_M(u^\star)=P_T.
\]
\end{lemma}

\begin{proof}
Puisque $p^\star$ appartient au relatif intérieur de la face $F$ et que la face active reste inchangée dans un voisinage de $u^\star$, la projection sur $M$ coïncide localement avec la projection orthogonale sur l'espace affine $\operatorname{aff}(F)$, suivie de la restriction au voisinage où le point projeté reste dans $F$.

Or la projection orthogonale sur un sous-espace affine est une application affine, donc différentiable partout, dont la différentielle est le projecteur orthogonal sur la direction linéaire de ce sous-espace. Cette direction est précisément $T_F$.

Ainsi,
\[
D\Pi_M(u^\star)=P_T.
\]
\end{proof}

\subsection{Lemme 2 : Hessien local de la distance carrée}

On considère la fonction
\[
\varphi(u)=d_M(u)^2=\|u-\Pi_M(u)\|^2.
\]

\begin{lemma}
Si $\Pi_M$ est différentiable en $u$, alors
\[
\nabla^2 \varphi(u)=2\bigl(I-D\Pi_M(u)\bigr).
\]
En particulier, sous les hypothèses locales précédentes,
\[
\nabla^2 \varphi(u^\star)=2(I-P_T)=2P_N.
\]
\end{lemma}

\begin{proof}
Le fait classique de l'analyse convexe est que
\[
\nabla\!\left(\frac12 d_M(u)^2\right)=u-\Pi_M(u).
\]
En multipliant par $2$, on obtient
\[
\nabla \varphi(u)=2(u-\Pi_M(u)).
\]
Si $\Pi_M$ est différentiable en $u$, on dérive cette identité :
\[
\nabla^2\varphi(u)=2I-2D\Pi_M(u)=2\bigl(I-D\Pi_M(u)\bigr).
\]
Au point $u^\star$, le lemme précédent donne $D\Pi_M(u^\star)=P_T$, donc
\[
\nabla^2\varphi(u^\star)=2(I-P_T)=2P_N.
\]
\end{proof}

\subsection{Structure spectrale de la projection locale}

\begin{proposition}
Sous les hypothèses locales précédentes, le spectre de $D\Pi_M(u^\star)$ est exactement
\[
\operatorname{Spec}(D\Pi_M(u^\star))
=
\{1\}^{\dim(T_F)}\cup \{0\}^{d-\dim(T_F)}.
\]
Par conséquent,
\[
\operatorname{Spec}(\nabla^2\varphi(u^\star))
=
\{0\}^{\dim(T_F)}\cup \{2\}^{d-\dim(T_F)}.
\]
\end{proposition}

\begin{proof}
Le lemme 1 montre que $D\Pi_M(u^\star)=P_T$. Or $P_T$ est un projecteur orthogonal, donc
\[
P_T^2=P_T,\qquad P_T^\top=P_T.
\]
Ses valeurs propres vérifient $\lambda^2=\lambda$, donc
\[
\lambda\in\{0,1\}.
\]
La multiplicité de la valeur propre $1$ est $\dim(T_F)$, et celle de $0$ est $d-\dim(T_F)$.

Comme
\[
\nabla^2\varphi(u^\star)=2(I-P_T),
\]
les valeurs propres correspondantes sont $0$ sur $T_F$ et $2$ sur $N_F$.
\end{proof}

\subsection{Fonctionnelle PCH locale}

On considère l'énergie PCH
\[
J(u;w,c)
=
\frac12\|u-w\|^2
+
\lambda_{\mathrm{ctx}}\,d_M(u)^2
+
\lambda_{\mathrm{ridge}}\|u\|^2.
\]

Sous les hypothèses locales ci-dessus, $J(\cdot;w,c)$ est de classe $C^2$ près de $u^\star$ et
\[
\nabla^2 J(u^\star;w,c)
=
I+2\lambda_{\mathrm{ridge}}I+2\lambda_{\mathrm{ctx}}(I-P_T).
\]
Autrement dit,
\[
\nabla^2 J(u^\star;w,c)
=
(1+2\lambda_{\mathrm{ridge}})I+2\lambda_{\mathrm{ctx}}(I-P_T).
\]

\begin{proof}
On dérive terme à terme :
\[
\nabla^2\Big(\frac12\|u-w\|^2\Big)=I,\qquad
\nabla^2\big(\lambda_{\mathrm{ridge}}\|u\|^2\big)=2\lambda_{\mathrm{ridge}}I,
\]
et, par le lemme précédent,
\[
\nabla^2\big(\lambda_{\mathrm{ctx}}d_M(u)^2\big)
=
2\lambda_{\mathrm{ctx}}(I-P_T).
\]
La somme donne la formule annoncée.
\end{proof}

\subsection{Jacobian local de l'opérateur proximal}

On définit l'opérateur proximal exact
\[
T_\eta(z;w,c)
=
\arg\min_u
\left\{
\frac12\|u-w\|^2
+
\lambda_{\mathrm{ctx}}\,d_M(u)^2
+
\lambda_{\mathrm{ridge}}\|u\|^2
+
\frac1{2\eta}\|u-z\|^2
\right\}.
\]

Soit $u^\star$ un point fixe local, c'est-à-dire
\[
u^\star=T_\eta(u^\star;w,c).
\]

\begin{theorem}[Loi résolvante locale exacte]
Sous les hypothèses précédentes,
\[
DT_\eta(u^\star;w,c)
=
\bigl(I+\eta \nabla^2J(u^\star;w,c)\bigr)^{-1}.
\]
En particulier, le spectre de $DT_\eta(u^\star;w,c)$ est donné par
\[
\operatorname{Spec}(DT_\eta(u^\star;w,c))
=
\left\{
\frac1{1+\eta(1+2\lambda_{\mathrm{ridge}})}
\right\}^{\dim(T_F)}
\cup
\left\{
\frac1{1+\eta(1+2\lambda_{\mathrm{ridge}}+2\lambda_{\mathrm{ctx}})}
\right\}^{d-\dim(T_F)}.
\]
\end{theorem}

\begin{proof}
Posons
\[
\Phi(u)=\frac12\|u-w\|^2+\lambda_{\mathrm{ctx}}\,d_M(u)^2+\lambda_{\mathrm{ridge}}\|u\|^2.
\]
L'opérateur proximal s'écrit
\[
T_\eta(z)=\arg\min_u\left\{\Phi(u)+\frac1{2\eta}\|u-z\|^2\right\}.
\]
La condition d'optimalité s'écrit
\[
\nabla\Phi(u)+\frac1\eta(u-z)=0.
\]
Au point fixe local $(u^\star,z^\star=u^\star)$, on applique le théorème des fonctions implicites. En différentiant par rapport à $z$, on obtient
\[
\left(\nabla^2\Phi(u^\star)+\frac1\eta I\right)DT_\eta(u^\star)
=
\frac1\eta I.
\]
Donc
\[
DT_\eta(u^\star)
=
\left(I+\eta\nabla^2\Phi(u^\star)\right)^{-1}.
\]
Or $\nabla^2\Phi(u^\star)=\nabla^2J(u^\star;w,c)$, ce qui donne la formule annoncée.

Enfin, comme
\[
\nabla^2J(u^\star;w,c)
=
(1+2\lambda_{\mathrm{ridge}})I+2\lambda_{\mathrm{ctx}}(I-P_T),
\]
les directions tangentes à la face active ont valeur propre
\[
\mu_T=1+2\lambda_{\mathrm{ridge}},
\]
et les directions normales ont valeur propre
\[
\mu_N=1+2\lambda_{\mathrm{ridge}}+2\lambda_{\mathrm{ctx}}.
\]
On obtient donc les valeurs propres de $DT_\eta(u^\star;w,c)$ en inversant $1+\eta\mu_T$ et $1+\eta\mu_N$.
\end{proof}

\subsection{Conséquence : constante de contraction observée}

Comme
\[
\mu_N\ge \mu_T,
\]
le rayon spectral est
\[
\rho\bigl(DT_\eta(u^\star;w,c)\bigr)
=
\frac1{1+\eta(1+2\lambda_{\mathrm{ridge}})}.
\]

Autrement dit, la constante de contraction locale observée ne dépend pas de la géométrie fine de la face active, mais seulement du terme quadratique injecté par
\[
\frac12\|u-w\|^2+\lambda_{\mathrm{ridge}}\|u\|^2,
\]
la contribution contextuelle n'affectant que les directions normales.

Ceci explique exactement les valeurs expérimentales :
\[
\lambda_{\mathrm{ridge}}=0.2 \quad\Longrightarrow\quad
\rho=\frac1{1+1.4}=\frac5{12}\approx 0.4166666667,
\]
\[
\lambda_{\mathrm{ridge}}=0.2,\ \eta=2 \quad\Longrightarrow\quad
\rho=\frac1{1+2\times 1.4}=\frac5{19}\approx 0.2631578947,
\]
\[
\lambda_{\mathrm{ridge}}=0.8 \quad\Longrightarrow\quad
\rho=\frac1{1+2.6}=\frac5{18}\approx 0.2777777778.
\]

\subsection{Ce qui est démontré et ce qui ne l'est pas}

Les résultats ci-dessus démontrent rigoureusement :

\begin{enumerate}
\item la formule locale exacte
\[
\nabla^2 d_M(u)^2=2(I-D\Pi_M(u));
\]

\item la structure spectrale de la projection convexe au voisinage d'une face active stable :
\[
\operatorname{Spec}(D\Pi_M(u^\star))\subset\{0,1\};
\]

\item la forme résolvante locale exacte du proximal :
\[
DT_\eta(u^\star;w,c)
=
\bigl(I+\eta\nabla^2J(u^\star;w,c)\bigr)^{-1};
\]

\item la loi locale de contraction :
\[
\rho
=
\frac1{1+\eta(1+2\lambda_{\mathrm{ridge}})}.
\]
\end{enumerate}

En revanche, la propriété suivante n'est pas démontrée par l'analyse locale seule :

\[
\text{l'universalité du même }H_{\mathrm{eff}}\text{ à travers tous les corpus, encodeurs et PCA.}
\]

Cette dernière affirmation requerrait un théorème supplémentaire de concentration géométrique des faces actives, ou une loi asymptotique de stabilité du projecteur tangent à travers la famille des hulls réels.

\subsection{Version Cas 3 local}

Sous les hypothèses locales de régularité de face active, on peut donc énoncer le Cas 3 local de manière rigoureuse :

\[
\boxed{
DT_\eta(u^\star;w,c)
=
\bigl(I+\eta H_{\mathrm{eff}}(c,w)\bigr)^{-1}
}
\]
avec
\[
H_{\mathrm{eff}}(c,w)
=
(1+2\lambda_{\mathrm{ridge}})I+2\lambda_{\mathrm{ctx}}(I-P_T).
\]

Dans ce sens précis, le noyau PCH possède bien une loi de type résolvant, et la constante de contraction observée n'est pas accidentelle : elle est la conséquence exacte de la structure convexe locale du problème.

\subsection{Remarque finale}

Si la face active est un singleton, alors $P_T=0$ et la formule devient
\[
DT_\eta(u^\star;w,c)
=
\frac1{1+\eta(1+2\lambda_{\mathrm{ridge}}+2\lambda_{\mathrm{ctx}})}I.
\]
Ce cas est particulier et ne décrit pas les runs expérimentaux consolidés, dans lesquels la direction tangentielle existe et gouverne le rayon spectral local.








%%%%%%%%%%%%%%%%%%%%%%%%%%%%%%%%%%%%%%%%%%%%%%%%%%%%%%%%%%%%%%%%%%%%%%%%%%%%%%%%%%%%%%%%%%%%%%%%%%%%%%%%%%%%%%%%%%%%%%%%%%%%%%%%%%%%%%%%%%%%%%%%%%%%%%%%%%%%%%%%%%%%%%%%%%%%%%%%%%%%%%%%%%%%%%%%%%%%%%%%%%%%%%%%%%%%%%%%%%%%%%%%%%%%%%%%%%%%%%%%%%%%%%%%%%%%%%%%%%%%%%%%%%%%%%%%%%%%%%%%%%%demonstration cas 3 universel :

\section{Cas 3 universel sous hypothèse de concentration géométrique}

\subsection{Cadre expérimental abstrait}

On considère une famille indexée par $\alpha$ de problèmes PCH, où $\alpha$ encode simultanément le corpus, l'encodeur, la réduction PCA, la dimension latente, la calibration, et les hyperparamètres.
Pour chaque $\alpha$, on dispose d'un ensemble admissible contextuel
\[
M_\alpha \subset \mathbb R^{d_\alpha},
\]
fermé, convexe et non vide, construit comme l'enveloppe convexe d'un ensemble fini de calibrations contextuelles.
On note
\[
\Pi_\alpha := \Pi_{M_\alpha}
\]
la projection euclidienne sur $M_\alpha$.

On considère l'énergie PCH locale
\[
J_\alpha(u;w,c)
=
\frac12\|u-w\|^2
+
\lambda_{\mathrm{ctx},\alpha}\,\dist(u,M_\alpha)^2
+
\lambda_{\mathrm{ridge},\alpha}\,\|u\|^2,
\]
et l'opérateur proximal associé
\[
T_{\eta,\alpha}(z;w,c)
=
\argmin_u
\left\{
\frac12\|u-w\|^2
+
\lambda_{\mathrm{ctx},\alpha}\,\dist(u,M_\alpha)^2
+
\lambda_{\mathrm{ridge},\alpha}\,\|u\|^2
+
\frac{1}{2\eta}\|u-z\|^2
\right\}.
\]

\subsection{Hypothèse supplémentaire de concentration géométrique}

\begin{assumption}[Concentration géométrique universelle des faces actives]
Il existe une constante $k_\star\in[0,1]$ et une matrice symétrique idempotente
\[
P_\star \in \mathbb R^{d_\alpha\times d_\alpha}
\]
telle que, pour toute famille admissible $\alpha$ et pour tout point fixe local $u_\alpha^\star$ du noyau PCH associé, on ait simultanément :

\begin{enumerate}
\item[(i)] $u_\alpha^\star$ appartient à un régime régulier où la projection est différentiable ;

\item[(ii)] la différentielle de la projection satisfait
\[
\|D\Pi_\alpha(u_\alpha^\star)-P_\star\|_{\mathrm{op}}
\le \varepsilon_\alpha,
\qquad \varepsilon_\alpha\to 0,
\]
uniformément sur la famille considérée ;

\item[(iii)] le projecteur limite est universel au sens où sa signature spectrale ne dépend ni du corpus, ni de l'encodeur, ni de la PCA, ni de la calibration, ni des hyperparamètres, sauf par des scalaires déjà explicités dans la fonctionnelle ;

\item[(iv)] les opérateurs locaux sont asymptotiquement isotropes sur les directions tangentielles et normales, de sorte que la courbure effective ne dépend asymptotiquement que d'un unique scalaire universel.
\end{enumerate}
\end{assumption}

Cette hypothèse est exactement le verrou manquant : elle remplace la simple observation empirique d'une loi de résolvant locale par une concentration uniforme des géométries de hulls.

\subsection{Lemme de base : Hessien de la distance carrée}

\begin{lemma}
Si $\Pi_\alpha$ est différentiable en $u$, alors
\[
\nabla^2 \dist(u,M_\alpha)^2
=
2\bigl(I-D\Pi_\alpha(u)\bigr).
\]
\end{lemma}

\begin{proof}
On utilise l'identité classique
\[
\nabla\left(\frac12 \dist(u,M_\alpha)^2\right)=u-\Pi_\alpha(u).
\]
En dérivant,
\[
\nabla^2\left(\frac12 \dist(u,M_\alpha)^2\right)
=
I-D\Pi_\alpha(u),
\]
d'où
\[
\nabla^2 \dist(u,M_\alpha)^2
=
2\bigl(I-D\Pi_\alpha(u)\bigr).
\]
\end{proof}

\subsection{Lemme de concentration du Hessien effectif}

\begin{lemma}
Sous l'hypothèse de concentration géométrique universelle, on a
\[
\nabla^2 \dist(u_\alpha^\star,M_\alpha)^2
\;\longrightarrow\;
2(I-P_\star)
\]
en norme opératorielle quand $\varepsilon_\alpha\to 0$.
\end{lemma}

\begin{proof}
Par le lemme précédent,
\[
\nabla^2 \dist(u_\alpha^\star,M_\alpha)^2
=
2\bigl(I-D\Pi_\alpha(u_\alpha^\star)\bigr).
\]
Comme
\[
\|D\Pi_\alpha(u_\alpha^\star)-P_\star\|_{\mathrm{op}}
\le \varepsilon_\alpha,
\]
on obtient
\[
\left\|
\nabla^2 \dist(u_\alpha^\star,M_\alpha)^2
-
2(I-P_\star)
\right\|_{\mathrm{op}}
\le 2\varepsilon_\alpha \to 0.
\]
\end{proof}

\subsection{Hessien universel effectif}

On pose
\[
H_\star
:=
(1+2\lambda_{\mathrm{ridge},\star})I
+
2\lambda_{\mathrm{ctx},\star}(I-P_\star),
\]
où $\lambda_{\mathrm{ridge},\star}$ et $\lambda_{\mathrm{ctx},\star}$ sont les constantes universelles imposées par la loi de concentration supposée.

\begin{theorem}[Cas 3 universel sous concentration géométrique]
Sous l'hypothèse précédente, le Hessien local du noyau PCH converge vers un opérateur universel
\[
\nabla^2 J_\alpha(u_\alpha^\star;w,c)\to H_\star,
\]
et le Jacobien local du proximal satisfait
\[
DT_{\eta,\alpha}(u_\alpha^\star;w,c)
\to
(I+\eta H_\star)^{-1}
\]
en norme opératorielle.

\end{theorem}

\begin{proof}
On a
\[
\nabla^2 J_\alpha(u_\alpha^\star;w,c)
=
I
+
2\lambda_{\mathrm{ridge},\alpha}I
+
2\lambda_{\mathrm{ctx},\alpha}\bigl(I-D\Pi_\alpha(u_\alpha^\star)\bigr).
\]
Par l'hypothèse de concentration géométrique,
\[
D\Pi_\alpha(u_\alpha^\star)\to P_\star,
\]
et, comme les coefficients sont supposés converger vers des constantes universelles,
\[
\lambda_{\mathrm{ridge},\alpha}\to \lambda_{\mathrm{ridge},\star},
\qquad
\lambda_{\mathrm{ctx},\alpha}\to \lambda_{\mathrm{ctx},\star}.
\]
Donc
\[
\nabla^2 J_\alpha(u_\alpha^\star;w,c)\to
(1+2\lambda_{\mathrm{ridge},\star})I
+
2\lambda_{\mathrm{ctx},\star}(I-P_\star)
=
H_\star.
\]

Par différentiation implicite de la condition d'optimalité du proximal,
\[
DT_{\eta,\alpha}(u_\alpha^\star;w,c)
=
\bigl(I+\eta \nabla^2 J_\alpha(u_\alpha^\star;w,c)\bigr)^{-1}.
\]
Le passage à la limite donne
\[
DT_{\eta,\alpha}(u_\alpha^\star;w,c)
\to
(I+\eta H_\star)^{-1}.
\]
\end{proof}

\subsection{Conséquence spectrale universelle}

Si $P_\star$ est un projecteur orthogonal de rang $r_\star$, alors
\[
\operatorname{Spec}(P_\star)=\{1\}^{r_\star}\cup\{0\}^{d-r_\star}.
\]
Donc
\[
\operatorname{Spec}(H_\star)
=
\left\{
1+2\lambda_{\mathrm{ridge},\star}
\right\}^{r_\star}
\cup
\left\{
1+2\lambda_{\mathrm{ridge},\star}+2\lambda_{\mathrm{ctx},\star}
\right\}^{d-r_\star}.
\]
Par conséquent,
\[
\operatorname{Spec}\bigl((I+\eta H_\star)^{-1}\bigr)
=
\left\{
\frac{1}{1+\eta(1+2\lambda_{\mathrm{ridge},\star})}
\right\}^{r_\star}
\cup
\left\{
\frac{1}{1+\eta(1+2\lambda_{\mathrm{ridge},\star}+2\lambda_{\mathrm{ctx},\star})}
\right\}^{d-r_\star}.
\]
Le rayon spectral est donc
\[
\rho_\star
=
\frac{1}{1+\eta(1+2\lambda_{\mathrm{ridge},\star})}.
\]

\subsection{Version universelle finale du Cas 3}

On peut alors formuler le Cas 3 universel ainsi :

\begin{theorem}[Cas 3 universel]
Sous l'hypothèse de concentration géométrique universelle des faces actives, la dynamique locale du noyau PCH est asymptotiquement gouvernée par un opérateur résolvant universel
\[
DT_{\eta,\alpha}(u_\alpha^\star;w,c)
\to
(I+\eta H_\star)^{-1},
\]
où $H_\star$ ne dépend ni du corpus, ni de l'encodeur, ni de la PCA, ni de la calibration, et où le rayon spectral local est universellement donné par
\[
\rho_\star
=
\frac{1}{1+\eta(1+2\lambda_{\mathrm{ridge},\star})}.
\]
\end{theorem}

\begin{proof}
C'est la réunion du lemme de concentration du Hessien effectif et du théorème précédent.
\end{proof}

\subsection{Remarque de portée}

La preuve ci-dessus est rigoureuse \emph{conditionnellement} à l'hypothèse de concentration géométrique.  
Sans cette hypothèse, on ne peut démontrer que la version locale du Cas 3.  
La version universelle est alors une conjecture de concentration structurelle des hulls réels.

\subsection{Lien avec le protocole expérimental}

Dans le protocole utilisé, la PCA est ajustée une seule fois sur les embeddings de calibration, puis réutilisée telle quelle sur les autres splits ; les contextes sont définis comme
\[
M_c=\operatorname{conv}\{E_{\theta_0}(x):x\in D_c\},
\]
où $D_c$ est l'ensemble de calibration du contexte $c$.  
Autrement dit, l'hypothèse de concentration porte exactement sur les objets géométriques réellement produits par le pipeline expérimental. 




%%%%%%%%%%%%%%%%%%%%%%%%%%%%%%%%%%%%%%%%%%%%%%%%%%%%%%%%%%%%%%%%%%%%%%%%%%%%%%%%%%%%%%%%%%%%%%%%%%%%%%%%%%%%%%%%%%%%%%%%%%%%%%%%%%%%%%%%%%%%%%%%%%%%%%%%%%%%%%%%%%%%%%%%%%%%%%%%%%%%%%%%%%%%%%%%%%%%%%%%%%%%%%%%%%%%%%%%%%%%%%%%%%%%%%%%%%%%%%%%%%%%%%%%%%%%%%%%%%%%%%%%%%%%%%%%%%%%%%%%%%%demonstration de la condition

\section{Hypothèse de concentration géométrique et réduction analytique du Cas 3 universel}

\subsection{Cadre local}

Soit \(M\subset \mathbb R^d\) un convexe compact non vide, et soit
\[
\Pi_M:\mathbb R^d\to M
\]
la projection euclidienne. On pose
\[
d_M(u)=\operatorname{dist}(u,M)=\|u-\Pi_M(u)\|.
\]
On considère la fonctionnelle PCH locale
\[
J(u;w,c)
=
\frac12\|u-w\|^2
+
\lambda_{\mathrm{ctx}}\, d_{M_c}(u)^2
+
\lambda_{\mathrm{ridge}}\|u\|^2.
\]

\subsection{Identité différentielle fondamentale}

\begin{lemma}[Hessien de la distance carrée]
Soit \(M\subset\mathbb R^d\) un convexe fermé.  
Si \(\Pi_M\) est différentiable en \(u\), alors
\[
\nabla^2 d_M(u)^2
=
2\bigl(I-D\Pi_M(u)\bigr).
\]
\end{lemma}

\begin{proof}
On utilise l'identité de Moreau pour l'enveloppe de la distance au carré :
\[
\nabla\!\left(\frac12 d_M(u)^2\right)=u-\Pi_M(u).
\]
En dérivant,
\[
\nabla^2\!\left(\frac12 d_M(u)^2\right)=I-D\Pi_M(u).
\]
En multipliant par \(2\), on obtient
\[
\nabla^2 d_M(u)^2
=
2\bigl(I-D\Pi_M(u)\bigr).
\]
\end{proof}

\subsection{Structure locale sur une face active}

Supposons maintenant que \(M=\operatorname{conv}(x_1,\dots,x_m)\) soit un polytope, et que
\[
p=\Pi_M(u^\star)
\]
appartienne au relatif intérieur d'une face active \(F\) de \(M\).

\begin{lemma}[Différentiabilité locale sur une face active]
Si la face active \(F\) reste inchangée dans un voisinage de \(u^\star\), alors
\[
D\Pi_M(u^\star)=P_F,
\]
où \(P_F\) est le projecteur orthogonal sur l'espace tangent \(T_F\) de la face.
\end{lemma}

\begin{proof}
Dans un voisinage où la face active reste fixe, la projection convexe sur \(M\) coïncide avec la projection affine orthogonale sur la face \(F\).  
La différentielle d'une projection orthogonale affine est exactement le projecteur orthogonal sur l'espace tangent.  
Donc \(D\Pi_M(u^\star)=P_F\).
\end{proof}

\begin{corollary}[Spectre local de la projection]
Comme \(P_F\) est un projecteur orthogonal,
\[
P_F^2=P_F,\qquad P_F^\top=P_F,
\]
et son spectre est contenu dans \(\{0,1\}\).  
Plus précisément,
\[
\operatorname{Spec}(D\Pi_M(u^\star))
=
\{1\}^{k}\cup\{0\}^{d-k},
\]
où \(k=\dim(T_F)\).
\end{corollary}

\subsection{Hypothèse de concentration géométrique}

On introduit maintenant l'hypothèse clé.

\begin{assumption}[Concentration géométrique uniforme]
Il existe un projecteur orthogonal déterministe \(P_0\) et une suite \(\varepsilon_n\to 0\) tels que, pour toute instance admissible
(corpus, encodeur, PCA, paramétrage proximal, contexte \(c\)), et pour tout point fixe \(u^\star\) du noyau PCH, on ait avec haute probabilité
\[
\|D\Pi_{M_c}(u^\star)-P_0\|_2\le \varepsilon_n.
\]
On suppose en outre que la dimension de la face active et les angles principaux du cône normal associé se concentrent de manière uniforme, de sorte que
\[
\operatorname{rank}(P_0)=k_{\mathrm{eff}}
\]
ne dépend pas des choix expérimentaux à l'ordre principal.
\end{assumption}

\subsection{Conséquence analytique pour le Hessien effectif}

\begin{proposition}[Hessien effectif sous concentration]
Sous l'hypothèse de concentration géométrique uniforme, on a au point fixe \(u^\star\),
\[
\nabla^2 d_{M_c}(u^\star)^2
=
2(I-P_0)+o(1).
\]
Par conséquent,
\[
\nabla^2 J(u^\star;w,c)
=
(1+2\lambda_{\mathrm{ridge}})I
+
2\lambda_{\mathrm{ctx}}(I-P_0)
+
o(1).
\]
\end{proposition}

\begin{proof}
Le premier énoncé suit immédiatement du lemme de Hessien :
\[
\nabla^2 d_{M_c}(u^\star)^2
=
2\bigl(I-D\Pi_{M_c}(u^\star)\bigr),
\]
et de l'hypothèse
\[
D\Pi_{M_c}(u^\star)=P_0+o(1).
\]
On en déduit
\[
\nabla^2 d_{M_c}(u^\star)^2
=
2(I-P_0)+o(1).
\]
En ajoutant les contributions quadratiques de
\[
\frac12\|u-w\|^2
\quad\text{et}\quad
\lambda_{\mathrm{ridge}}\|u\|^2,
\]
on obtient
\[
\nabla^2 J(u^\star;w,c)
=
I+2\lambda_{\mathrm{ridge}}I
+2\lambda_{\mathrm{ctx}}(I-P_0)
+o(1),
\]
soit
\[
\nabla^2 J(u^\star;w,c)
=
(1+2\lambda_{\mathrm{ridge}})I
+
2\lambda_{\mathrm{ctx}}(I-P_0)
+
o(1).
\]
\end{proof}

\subsection{Linéarisation du prox}

Soit
\[
T_\eta^{\mathrm{PCH}}(z;w,c)
=
\arg\min_u
\left\{
\frac12\|u-w\|^2
+
\lambda_{\mathrm{ctx}} d_{M_c}(u)^2
+
\lambda_{\mathrm{ridge}}\|u\|^2
+
\frac{1}{2\eta}\|u-z\|^2
\right\}.
\]

\begin{theorem}[Forme résolvante locale]
Sous l'hypothèse de concentration géométrique uniforme, le jacobien au point fixe vérifie
\[
DT_\eta^{\mathrm{PCH}}(u^\star)
=
\Bigl(I+\eta H_{\mathrm{eff}}\Bigr)^{-1}
+
o(1),
\]
où
\[
H_{\mathrm{eff}}
=
(1+2\lambda_{\mathrm{ridge}})I
+
2\lambda_{\mathrm{ctx}}(I-P_0).
\]
\end{theorem}

\begin{proof}
Dans le régime local où \(J\) est différentiable et strictement convexe, le prox résout l'équation stationnaire
\[
0
=
\nabla J(u^\star;w,c)+\frac1\eta(u^\star-z).
\]
Par le théorème des fonctions implicites appliqué à l'opérateur de premier ordre, la différentielle de la solution par rapport à \(z\) est le résolvant du Hessien local :
\[
DT_\eta^{\mathrm{PCH}}(u^\star)
=
\bigl(I+\eta \nabla^2 J(u^\star;w,c)\bigr)^{-1}.
\]
En remplaçant \(\nabla^2 J(u^\star;w,c)\) par son expression effective précédente, on obtient
\[
DT_\eta^{\mathrm{PCH}}(u^\star)
=
\Bigl(I+\eta H_{\mathrm{eff}}\Bigr)^{-1}
+
o(1).
\]
\end{proof}

\subsection{Cas 3 universel}

\begin{theorem}[Cas 3 universel conditionnel à la concentration]
Si l'hypothèse de concentration géométrique uniforme est vraie avec un projecteur limite \(P_0\) indépendant du corpus, de l'encodeur, de la PCA et des hyperparamètres, alors le Cas 3 universel est vrai au premier ordre :
\[
DT_\eta^{\mathrm{PCH}}(u^\star)
=
\Bigl(I+\eta H_{\mathrm{eff}}\Bigr)^{-1}
+
o(1),
\]
avec \(H_{\mathrm{eff}}\) universel à l'ordre principal.
\end{theorem}

\begin{proof}
Par la proposition précédente, la seule dépendance non universelle de la différentielle locale du prox est portée par \(P_0\).  
Si \(P_0\) est lui-même universel, alors \(H_{\mathrm{eff}}\) est universel à l'ordre principal.  
La conclusion suit.
\end{proof}

\subsection{Ce qu'il reste à prouver pour fermer complètement le Cas 3}

La preuve ci-dessus réduit le problème à une seule question probabiliste :
\[
\|D\Pi_{M_c}(u^\star)-P_0\|_2\to 0
\quad\text{uniformément sur la famille admissible.}
\]

Autrement dit, il reste à démontrer un théorème de concentration pour les faces actives des hulls convexes construits à partir des embeddings de calibration.

Sans modèle génératif explicite pour ces embeddings, cette étape ne peut pas être fermée uniquement par des arguments analytiques internes au noyau PCH.




%%%%%%%%%%%%%%%%%%%%%%%%%%%%%%%%%%%%%%%%%%%%%%%%%%%%%%%%%%%%%%%%%%%%%%%%%%%%%%%%%%%%%%%%%%%%%%%%%%%%%%%%%%%%%%%%%%%%%%%%%%%%%%%%%%%%%%%%%%%%%%%%%%%%%%%%%%%%%%%%%%%%%%%%%%%%%%%%%%%%%%%%%%%%%%%%%%%%%%%%%%%%%%%%%%%%%%%%%%%%%%%%%%%%%%%%%%%%%%%%%%%%%%%%%%%%%%%%%%%%%%%%%%%%%%%%%%%%%%%%%%%model proba



\section{Modèle probabiliste explicite et théorème de concentration}

\subsection{Pourquoi un modèle probabiliste est nécessaire}

Si les embeddings de calibration sont considérés comme des objets entièrement déterministes,
alors aucune phrase du type
\[
\|D\Pi_{M_c}(u^\star)-P_0\|_2\to 0
\]
ne peut être démontrée probabilistiquement.

Pour obtenir un théorème de concentration, il faut donc postuler un modèle aléatoire pour les embeddings de calibration, puis montrer que, sous ce modèle, les projecteurs de projection convexe et les Jacobiens du prox se concentrent.

\subsection{Modèle génératif des embeddings de calibration}

On fixe un entier \(k\ll d\), une matrice orthonormale \(V\in\mathbb R^{d\times k}\),
et un convexe compact \(K\subset\mathbb R^k\).

On suppose que les embeddings de calibration d'un contexte \(c\) sont i.i.d. de la forme
\[
X_i=\mu_c+V S_i+\sigma \xi_i,
\qquad i=1,\dots,n,
\]
où :

\begin{itemize}
    \item \(\mu_c\in\mathbb R^d\) est le barycentre contextuel ;
    \item \(S_i\) sont i.i.d., supportés dans \(K\), de densité bornée inférieurement et supérieurement sur \(K\) ;
    \item \(\xi_i\) sont i.i.d. isotropes, centrés, \(\psi_2\)-sous-gaussiens, indépendants des \(S_i\) ;
    \item \(\sigma\ge 0\) contrôle le bruit transversal ;
    \item \(V\) représente la géométrie intrinsèque du contexte.
\end{itemize}

Le hull empirique est
\[
M_n=\operatorname{conv}(X_1,\dots,X_n),
\]
et le hull population est
\[
M_\star=\mu_c+VK.
\]

\subsection{Hypothèse géométrique de non-dégénérescence}

Soit \(u^\star\) un point fixe du noyau PCH, et
\[
p^\star=\Pi_{M_\star}(u^\star).
\]
On suppose que \(p^\star\) appartient au relatif intérieur d'une face active \(F_\star\) de \(M_\star\),
et que cette face est \emph{stable}, au sens suivant :

\[
\operatorname{dist}\!\left(p^\star,\partial F_\star\right)\ge \tau>0,
\]
et l'angle dièdre minimal entre \(F_\star\) et toute face concurrente est borné inférieurement par \(\alpha>0\).

Cette hypothèse est la version précise de la \emph{concentration géométrique} : elle dit que le point projeté ne touche pas une zone singulière du hull.

\subsection{Concentration du hull empirique}

On introduit la fonction de support
\[
h_A(v)=\sup_{x\in A}\langle v,x\rangle.
\]

\begin{theorem}[Concentration uniforme du support]
Pour tout \(\delta\in(0,1)\), il existe une constante \(C>0\), dépendant seulement des paramètres sous-gaussiens et de la géométrie de \(K\), telle que, avec probabilité au moins \(1-\delta\),
\[
\sup_{\|v\|=1}\bigl|h_{M_n}(v)-h_{M_\star}(v)\bigr|
\le
C\left(
\sqrt{\frac{k+\log(1/\delta)}{n}}
+
\sigma\sqrt{\frac{d+\log(1/\delta)}{n}}
\right).
\]
\end{theorem}

\begin{proof}[Preuve]
On écrit
\[
h_{M_n}(v)=\max_{1\le i\le n}\langle v,X_i\rangle.
\]
Comme \(X_i=\mu_c+VS_i+\sigma\xi_i\), on a
\[
\langle v,X_i\rangle
=
\langle v,\mu_c\rangle
+
\langle V^\top v,S_i\rangle
+
\sigma\langle v,\xi_i\rangle.
\]
Le premier terme est déterministe, le second se concentre par borne uniforme sur une famille de variables sous-gaussiennes indexée par la sphère, et le troisième se contrôle par inégalités standards de concentration sous-gaussienne et union bound sur un \(\varepsilon\)-net de \(\mathbb S^{d-1}\).

En combinant :
\begin{itemize}
    \item la concentration du supremum empirique,
    \item le contrôle par \(\varepsilon\)-net,
    \item l'invariance de la norme opérateur sous rotation,
\end{itemize}
on obtient l'estimation annoncée.
\end{proof}

\subsection{Convergence des projections}

\begin{theorem}[Stabilité de la projection convexe]
Sous l'hypothèse précédente, il existe une constante \(C'>0\) telle que, avec probabilité au moins \(1-\delta\),
\[
\sup_{\|u-u^\star\|\le r}
\|\Pi_{M_n}(u)-\Pi_{M_\star}(u)\|
\le
C'
\left(
\sqrt{\frac{k+\log(1/\delta)}{n}}
+
\sigma\sqrt{\frac{d+\log(1/\delta)}{n}}
\right),
\]
pour tout rayon \(r\) contenu dans le tube normal où \(\Pi_{M_\star}\) est \(C^1\).
\end{theorem}

\begin{proof}
Pour les convexes fermés, la projection euclidienne est 1-Lipschitz par rapport à l'argument et stable sous convergence de Hausdorff des ensembles convexes.

Or la borne précédente sur les fonctions de support implique une borne de Hausdorff :
\[
d_H(M_n,M_\star)
\le
\sup_{\|v\|=1}\bigl|h_{M_n}(v)-h_{M_\star}(v)\bigr|.
\]
Le résultat suit par stabilité de la projection sous perturbation de Hausdorff, dans le tube normal où la projection reste unique et régulière.
\end{proof}

\subsection{Concentration du Jacobien de projection}

\begin{theorem}[Concentration du différentiel de projection]
Sous les hypothèses de stabilité de face et de concentration précédente, il existe un projecteur orthogonal \(P_0\) tel que, avec probabilité au moins \(1-\delta\),
\[
\|D\Pi_{M_n}(u^\star)-P_0\|_2
\le
C''
\left(
\sqrt{\frac{k+\log(1/\delta)}{n}}
+
\sigma\sqrt{\frac{d+\log(1/\delta)}{n}}
\right).
\]
\end{theorem}

\begin{proof}
Par hypothèse, la face active de \(M_\star\) au point \(p^\star\) est stable et de rang tangent fixé.
Donc
\[
D\Pi_{M_\star}(u^\star)=P_0.
\]
La stabilité en Hausdorff du hull entraîne la stabilité de la face active tant que la perturbation est inférieure à la marge \(\tau\) et que l'angle minimal \(\alpha\) reste non nul.

Dans ce régime, le différentiel de projection dépend Lipschitz-continuement de la géométrie locale du bord convexe.
On obtient donc
\[
\|D\Pi_{M_n}(u^\star)-D\Pi_{M_\star}(u^\star)\|_2
\le
C''\, d_H(M_n,M_\star),
\]
et le résultat suit de l'estimation précédente.
\end{proof}

\subsection{Théorème de concentration pour le Hessien effectif}

\begin{corollary}[Concentration du Hessien effectif]
Sous les hypothèses précédentes,
\[
\nabla^2 d_{M_n}(u^\star)^2
=
2(I-P_0)+o_{\mathbb P}(1),
\]
et donc
\[
\nabla^2 J(u^\star;w,c)
=
(1+2\lambda_{\mathrm{ridge}})I
+
2\lambda_{\mathrm{ctx}}(I-P_0)
+
o_{\mathbb P}(1).
\]
\end{corollary}

\begin{proof}
C'est une conséquence directe de l'identité
\[
\nabla^2 d_{M_n}(u^\star)^2=2(I-D\Pi_{M_n}(u^\star))
\]
et du théorème de concentration du différentiel de projection.
\end{proof}

\subsection{Conséquence pour le Cas 3 universel}

\begin{theorem}[Cas 3 universel sous concentration]
Si les embeddings de calibration satisfont le modèle probabiliste ci-dessus,
et si les paramètres géométriques \((k,\tau,\alpha,\sigma)\) sont uniformément contrôlés sur la famille de corpus, encodeurs, PCA et hyperparamètres considérée,
alors
\[
DT_\eta^{\mathrm{PCH}}(u^\star)
=
\Bigl(I+\eta H_{\mathrm{eff}}\Bigr)^{-1}
+
o_{\mathbb P}(1),
\]
avec
\[
H_{\mathrm{eff}}
=
(1+2\lambda_{\mathrm{ridge}})I
+
2\lambda_{\mathrm{ctx}}(I-P_0).
\]
\end{theorem}

\begin{proof}
Par différentiation implicite de l'équation stationnaire du prox,
\[
0=\nabla J(u^\star;w,c)+\frac1\eta(u^\star-z),
\]
on a localement
\[
DT_\eta^{\mathrm{PCH}}(u^\star)
=
\bigl(I+\eta\nabla^2J(u^\star;w,c)\bigr)^{-1}.
\]
Le corollaire précédent fournit la concentration de \(\nabla^2J(u^\star;w,c)\) autour de \(H_{\mathrm{eff}}\).  
On en déduit la formule annoncée.
\end{proof}

\subsection{Ce qu'il reste à vérifier pour transformer ce résultat en preuve complète}

Le point vraiment dur n'est plus PCH lui-même, mais le modèle probabiliste des embeddings.  
Pour fermer complètement le Cas 3 universel, il faut vérifier empiriquement ou démontrer analytiquement que :

\begin{enumerate}
    \item la famille des embeddings de calibration admet un modèle sous-gaussien ou à support convexe contrôlé ;
    \item la face active du hull population est stable avec marge \(\tau>0\) ;
    \item les paramètres \((k,\tau,\alpha,\sigma)\) sont uniformément bornés sur la classe de corpus et d'encodeurs étudiée ;
    \item la PCA ne détruit pas la séparation géométrique utile ;
    \item les fluctuations dues au routage ne déplacent pas le point fixe hors du tube normal stable.
\end{enumerate}

Sous ces conditions, l'hypothèse de concentration géométrique devient un théorème de concentration standard sur les convexes aléatoires, et le Cas 3 universel suit.




%%%%%%%%%%%%%%%%%%%%%%%%%%%%%%%%%%%%%%%%%%%%%%%%%%%%%%%%%%%%%%%%%%%%%%%%%%%%%%%%%%%%%%%%%%%%%%%%%%%%%%%%%%%%%%%%%%%%%%%%%%%%%%%%%%%%%%%%%%%%%%%%%%%%%%%%%%%%%%%%%%%%%%%%%%%%%%%%%%%%%%%%%%%%%%%%%%%%%%%%%%%%%%%%%%%%%%%%%%%%%%%%%%%%%%%%%%%%%%%%%%%%%%%%%%%%%%%%%%%%%%%%%%%%%%%%%%%%%%%%%%%


\section{Concentration du différentiel de projection autour d'un projecteur limite}

\subsection{Modèle probabiliste explicite}

Soit \(c\) un contexte fixé. On considère les embeddings de calibration
\[
X_1,\dots,X_n \in \mathbb R^d
\]
comme des variables aléatoires i.i.d. de la forme
\[
X_i=\mu_c+V S_i+\sigma \xi_i,
\]
où :

\begin{itemize}
    \item \(\mu_c\in\mathbb R^d\) est le barycentre contextuel ;
    \item \(V\in\mathbb R^{d\times k}\) a colonnes orthonormées et représente la géométrie intrinsèque du contexte ;
    \item \(S_i\in\mathbb R^k\) sont i.i.d., supportées sur un convexe compact \(K\subset\mathbb R^k\), avec densité bornée supérieurement et inférieurement sur \(K\) ;
    \item \(\xi_i\in\mathbb R^d\) sont i.i.d., centrées, isotropes et \(\psi_2\)-sous-gaussiennes ;
    \item \(\sigma\ge 0\) contrôle l'épaisseur transverse du nuage.
\end{itemize}

On définit le hull empirique
\[
M_n=\conv(X_1,\dots,X_n),
\]
et le hull population
\[
M_\star=\mu_c+VK.
\]

\subsection{Hypothèses géométriques}

Soit \(u^\star\in\mathbb R^d\) le point fixe pertinent du noyau PCH, et
\[
p_\star=\Pi_{M_\star}(u^\star).
\]
On suppose que \(p_\star\) appartient au relatif intérieur d'une face \(F_\star\) de \(M_\star\), et que cette face est \emph{stable} au sens suivant :

\begin{enumerate}
    \item \(\operatorname{dist}(p_\star,\partial F_\star)\ge \tau>0\) ;
    \item le cône normal de \(F_\star\) est isolé des autres cônes normaux par une marge angulaire \(\alpha>0\).
\end{enumerate}

On pose
\[
P_0 := D\Pi_{M_\star}(u^\star).
\]
Sous l'hypothèse de stabilité de face, \(P_0\) est le projecteur orthogonal sur l'espace tangent de \(F_\star\).

\subsection{Lemme de concentration du support}

Pour un convexe compact \(A\subset\mathbb R^d\), on note sa fonction de support
\[
h_A(v)=\sup_{x\in A}\langle v,x\rangle.
\]

\begin{lemma}[Concentration uniforme du support]
Pour tout \(\delta\in(0,1)\), il existe une constante \(C>0\), dépendant uniquement des bornes sous-gaussiennes et de la géométrie de \(K\), telle que
\[
\mathbb P\!\left(
\sup_{\|v\|=1}\bigl|h_{M_n}(v)-h_{M_\star}(v)\bigr|
\le
C\Bigl(
\sqrt{\tfrac{k+\log(1/\delta)}{n}}
+
\sigma\sqrt{\tfrac{d+\log(1/\delta)}{n}}
\Bigr)
\right)\ge 1-\delta.
\]
\end{lemma}

\begin{proof}
On combine une borne de concentration pour les maxima empiriques de variables sous-gaussiennes avec un argument de \(\varepsilon\)-net sur \(\mathbb S^{d-1}\). Le terme intrinsèque en \(k\) provient de la partie \(V S_i\), tandis que le terme en \(d\) provient du bruit transverse \(\sigma\xi_i\). Le contrôle uniforme suit par union bound sur le filet sphérique et passage à la limite.
\end{proof}

\subsection{Contrôle de la distance de Hausdorff}

Pour des convexes compacts,
\[
d_H(A,B)=\sup_{\|v\|=1}|h_A(v)-h_B(v)|.
\]
Ainsi, avec le lemme précédent,
\[
d_H(M_n,M_\star)
=
O_{\mathbb P}\!\left(
\sqrt{\frac{k}{n}}
+
\sigma\sqrt{\frac{d}{n}}
\right).
\]

\subsection{Stabilité de la face active}

\begin{lemma}[Persistance de la face active]
Il existe une constante \(c_0>0\) ne dépendant que de \(\tau\) et \(\alpha\) telle que, si
\[
d_H(M_n,M_\star)\le c_0,
\]
alors la projection de \(u^\star\) sur \(M_n\) tombe dans la face correspondante \(F_n\), combinatoirement équivalente à \(F_\star\), et
\[
\Pi_{M_n}
\quad\text{est différentiable en }u^\star.
\]
\end{lemma}

\begin{proof}
La marge \(\tau\) empêche la projection limite d'approcher une arête singulière ou une intersection de faces concurrentes. La marge angulaire \(\alpha\) assure l'isolement du cône normal de \(F_\star\). Par stabilité des faces exposées sous perturbation de Hausdorff petite devant ces marges, le même ensemble d'indices actifs reste optimal pour \(M_n\). La projection reste alors localement régulière.
\end{proof}

\subsection{Forme du différentiel de projection}

\begin{proposition}[Différentiel local]
Sous la stabilité de face,
\[
D\Pi_{M_n}(u^\star)=P_{T_n},
\]
où \(P_{T_n}\) est le projecteur orthogonal sur l'espace tangent de la face active \(F_n\).
De plus,
\[
D\Pi_{M_\star}(u^\star)=P_0=P_{T_\star}.
\]
\end{proposition}

\begin{proof}
Sur l'intérieur relatif d'une face active fixe, la projection convexe est affine le long des directions tangentes et annihile les directions normales. Sa différentielle est donc exactement le projecteur orthogonal sur l'espace tangent.
\end{proof}

\subsection{Concentration du projecteur différentiel}

\begin{theorem}[Concentration de \(D\Pi_{M_n}(u^\star)\) autour de \(P_0\)]
Sous le modèle probabiliste ci-dessus et sous les hypothèses de marge \(\tau,\alpha>0\), il existe une constante \(C'>0\) telle que, pour tout \(\delta\in(0,1)\),
\[
\mathbb P\!\left(
\bigl\|D\Pi_{M_n}(u^\star)-P_0\bigr\|_2
\le
C'\Bigl(
\sqrt{\tfrac{k+\log(1/\delta)}{n}}
+
\sigma\sqrt{\tfrac{d+\log(1/\delta)}{n}}
\Bigr)
\right)\ge 1-\delta.
\]
\end{theorem}

\begin{proof}
Par le lemme de concentration du support, on a avec probabilité au moins \(1-\delta\),
\[
d_H(M_n,M_\star)\le \varepsilon_n,
\qquad
\varepsilon_n=
C\Bigl(
\sqrt{\tfrac{k+\log(1/\delta)}{n}}
+
\sigma\sqrt{\tfrac{d+\log(1/\delta)}{n}}
\Bigr).
\]
Si \(n\) est assez grand pour que \(\varepsilon_n\le c_0\), alors la face active persiste et
\[
D\Pi_{M_n}(u^\star)=P_{T_n},\qquad D\Pi_{M_\star}(u^\star)=P_0.
\]
La stabilité des sous-espaces actifs sous perturbation de Hausdorff implique alors, par un argument de perturbation de sous-espaces de type Davis--Kahan,
\[
\|P_{T_n}-P_0\|_2\le C'\,\varepsilon_n.
\]
En combinant les deux inégalités, on obtient le résultat annoncé.
\end{proof}

\subsection{Corollaire : concentration du Hessien effectif}

Comme
\[
\nabla^2 d_{M_n}(u^\star)^2
=
2\bigl(I-D\Pi_{M_n}(u^\star)\bigr),
\]
on déduit immédiatement
\[
\nabla^2 d_{M_n}(u^\star)^2
=
2(I-P_0)+o_{\mathbb P}(1).
\]
Par conséquent, pour l'énergie PCH,
\[
\nabla^2 J_{\mathrm{PCH}}(u^\star;w,c)
=
(1+2\lambda_{\mathrm{ridge}})I
+
2\lambda_{\mathrm{ctx}}(I-P_0)
+
o_{\mathbb P}(1).
\]

\subsection{Cas 3 universel}

Si, en plus, la famille entière des corpus, encodeurs, réductions PCA et hyperparamètres considérés satisfait une \emph{universality hypothesis} au sens où le projecteur limite \(P_0\) est le même à travers la famille après normalisation, alors
\[
D\Pi_{M_n}(u^\star)\xrightarrow{\mathbb P} P_0
\]
uniformément sur cette famille, et
\[
DT_\eta^{\mathrm{PCH}}(u^\star)
=
\bigl(I+\eta H_{\mathrm{eff}}\bigr)^{-1}
+
o_{\mathbb P}(1),
\]
avec
\[
H_{\mathrm{eff}}
=
(1+2\lambda_{\mathrm{ridge}})I
+
2\lambda_{\mathrm{ctx}}(I-P_0).
\]

\subsection{Lecture conceptuelle}

La bonne lecture est maintenant nette :

\begin{itemize}
    \item la concentration du projecteur n'est pas un fait magique du prox ;
    \item elle découle d'abord d'une concentration géométrique du hull empirique ;
    \item la stabilité des faces actives transmet ensuite cette concentration au différentiel de projection ;
    \item le Hessien effectif n'est qu'une conséquence de cette stabilité géométrique.
\end{itemize}



%%%%%%%%%%%%%%%%%%%%%%%%%%%%%%%%%%%%%%%%%%%%%%%%%%%%%%%%%%%%%%%%%%%%%%%%%%%%%%%%%%%%%%%%%%%%%%%%%%%%%%%%%%%%%%%%%%%%%%%%%%%%%%%%%%%%%%%%%%%%%%%%%%%%%%%%%%%%%%%%%%%%%%%%%%%%%%%%%%%%%%%%%%%%%%%%%%%%%%%%%%%%%%%%%%%%%%%%%%%%%%%%%%%%%%%%%%%%%%%%%%%%%%%%%%%%%%%%%%%%%%%%%%%%%%%%%%%%%%%%%%%suite 

\subsection{Concentration du différentiel de projection autour d'un projecteur limite}
\label{subsec:projection_diff_concentration}

\paragraph{But.}
On veut montrer que, sous un modèle probabiliste explicite et sous une hypothèse de stabilité géométrique du face active, le différentiel de projection sur le hull empirique converge vers un projecteur limite universel :
\[
D\Pi_{M_n}(u^\star)\;\longrightarrow\;P_0,
\qquad
P_0:=D\Pi_{M_\star}(u^\star).
\]

\subsubsection{Modèle probabiliste et objets géométriques}

Soit \(c\) un contexte fixé. On considère les embeddings de calibration
\[
X_1,\dots,X_n\in\mathbb R^d
\]
comme i.i.d. selon le modèle
\[
X_i=\mu_c+VS_i+\sigma\xi_i,
\]
où :
\begin{itemize}
    \item \(\mu_c\in\mathbb R^d\) est le barycentre contextuel ;
    \item \(V\in\mathbb R^{d\times k}\) a colonnes orthonormées ;
    \item \(S_i\in\mathbb R^k\) sont i.i.d., supportées dans un convexe compact \(K\subset\mathbb R^k\), avec densité bornée inférieurement et supérieurement sur \(K\) ;
    \item \(\xi_i\in\mathbb R^d\) sont i.i.d., centrées, isotropes et \(\psi_2\)-sous-gaussiennes ;
    \item \(\sigma\ge 0\) contrôle le bruit transverse.
\end{itemize}

On définit le hull empirique
\[
M_n:=\operatorname{conv}(X_1,\dots,X_n),
\]
et le hull population
\[
M_\star:=\mu_c+VK.
\]

Soit \(u^\star\in\mathbb R^d\) le point fixe pertinent du noyau PCH, et
\[
p_\star:=\Pi_{M_\star}(u^\star).
\]
On suppose que \(p_\star\) appartient au relatif intérieur d'une face \(F_\star\) de \(M_\star\).

\begin{assumption}[Séparation géométrique]
Il existe \(\tau>0\) et \(\alpha>0\) tels que :
\begin{enumerate}
    \item \(\operatorname{dist}(p_\star,\partial F_\star)\ge \tau\) ;
    \item le cône normal \(N_{M_\star}(p_\star)\) est isolé des autres cônes normaux par une marge angulaire \(\alpha\).
\end{enumerate}
\end{assumption}

On pose
\[
P_0:=D\Pi_{M_\star}(u^\star).
\]
Sous l'hypothèse de stabilité de face, \(P_0\) est le projecteur orthogonal sur l'espace tangent de \(F_\star\).

\subsubsection{Étape 1 : concentration de la fonction de support}

Pour un convexe compact \(A\subset\mathbb R^d\), on note sa fonction de support
\[
h_A(v)=\sup_{x\in A}\langle v,x\rangle,
\qquad \|v\|=1.
\]

\begin{lemma}[Concentration uniforme du support]
Pour tout \(\delta\in(0,1)\), il existe une constante \(C>0\), dépendant uniquement des bornes sous-gaussiennes et de la géométrie de \(K\), telle que
\[
\mathbb P\!\left(
\sup_{\|v\|=1}\bigl|h_{M_n}(v)-h_{M_\star}(v)\bigr|
\le
C\left(
\sqrt{\frac{k+\log(1/\delta)}{n}}
+
\sigma\sqrt{\frac{d+\log(1/\delta)}{n}}
\right)
\right)\ge 1-\delta.
\]
\end{lemma}

\begin{proof}
Fixons \(v\in\mathbb S^{d-1}\). Par définition,
\[
h_{M_n}(v)=\max_{1\le i\le n}\langle v,X_i\rangle.
\]
Or \(\langle v,X_i\rangle=\langle v,\mu_c\rangle+\langle V^\top v,S_i\rangle+\sigma\langle v,\xi_i\rangle\).
La partie intrinsèque vit dans \(\mathbb R^k\), donc la concentration de son maximum dépend de \(k\), tandis que le bruit transverse produit un terme sous-gaussien de dimension ambiante \(d\).

On applique ensuite un argument de filet \(\varepsilon\)-net sur \(\mathbb S^{d-1}\), puis un passage à la limite par continuité lipschitzienne de la fonction support en \(v\). On obtient la borne uniforme annoncée.
\end{proof}

\begin{corollary}[Hausdorff]
Avec la même probabilité,
\[
d_H(M_n,M_\star)
=
\sup_{\|v\|=1}\bigl|h_{M_n}(v)-h_{M_\star}(v)\bigr|
\le
C\left(
\sqrt{\frac{k+\log(1/\delta)}{n}}
+
\sigma\sqrt{\frac{d+\log(1/\delta)}{n}}
\right).
\]
En particulier,
\[
d_H(M_n,M_\star)=O_{\mathbb P}\!\left(\sqrt{\frac{k}{n}}+\sigma\sqrt{\frac{d}{n}}\right).
\]
\end{corollary}

\subsubsection{Étape 2 : stabilité de la face active}

\begin{lemma}[Persistance de la face active]
Il existe une constante \(c_0>0\), ne dépendant que de \(\tau\) et \(\alpha\), telle que si
\[
d_H(M_n,M_\star)\le c_0,
\]
alors :
\begin{enumerate}
    \item la projection \(p_n:=\Pi_{M_n}(u^\star)\) appartient à une face \(F_n\) combinatoirement équivalente à \(F_\star\) ;
    \item la projection \(\Pi_{M_n}\) est différentiable en \(u^\star\) ;
    \item l'ensemble actif de contraintes reste constant dans un voisinage de \(u^\star\).
\end{enumerate}
\end{lemma}

\begin{proof}
La marge \(\tau\) assure que \(p_\star\) n'est pas près du bord de la face \(F_\star\). La marge angulaire \(\alpha\) empêche la collision avec un autre cône normal concurrent. Il existe donc un voisinage de \(M_\star\) en distance de Hausdorff dans lequel la face exposée par \(u^\star\) reste la même à homéomorphisme combinatoire près.

Comme les polytopes sont des convexes définis par un nombre fini de faces, la stabilité de la face active sous petite perturbation de Hausdorff est une conséquence standard de la non-dégénérescence : les identifiants actifs ne changent pas, et la projection reste localement affine le long des directions tangentes à la face active.
\end{proof}

\subsubsection{Étape 3 : forme du différentiel de projection}

\begin{proposition}[Différentiel local de la projection]
Sous la stabilité de la face active,
\[
D\Pi_{M_n}(u^\star)=P_{T_n},
\]
où \(P_{T_n}\) est le projecteur orthogonal sur l'espace tangent \(T_n:=\operatorname{Tan}(F_n)\). De même,
\[
D\Pi_{M_\star}(u^\star)=P_0=P_{T_\star}.
\]
\end{proposition}

\begin{proof}
Sur le relatif intérieur d'une face active fixe, la projection euclidienne sur un convexe polyédral est affine en première approximation : elle conserve les directions tangentes à la face et annule les directions normales. Sa différentielle est donc exactement le projecteur orthogonal sur l'espace tangent de la face active.
\end{proof}

\subsubsection{Étape 4 : contrôle de l'écart entre les projecteurs tangents}

\begin{lemma}[Perturbation du projecteur tangent]
Il existe une constante \(C_1>0\), dépendant uniquement du conditionnement géométrique de \(F_\star\) (en particulier de sa marge \(\tau\), de son angle d'ouverture \(\alpha\) et d'un minorant de singularité du système de directions actives), telle que, pour \(n\) assez grand,
\[
\|P_{T_n}-P_0\|_2 \le C_1\, d_H(M_n,M_\star).
\]
\end{lemma}

\begin{proof}
Soient \(u_1,\dots,u_m\) des directions indépendantes engendrant \(T_\star\), et \(U_\star=[u_1\,\cdots\,u_m]\). De même, \(U_n=[u_1^{(n)}\,\cdots\,u_m^{(n)}]\) engendre \(T_n\), avec
\[
\|U_n-U_\star\|_2\le C_2\, d_H(M_n,M_\star)
\]
pour une constante \(C_2\) ne dépendant que de la géométrie locale de la face.

Le projecteur orthogonal sur \(\operatorname{Im}(U)\) s'écrit
\[
P_U=U(U^\top U)^{-1}U^\top.
\]
Si \(\sigma_{\min}(U_\star)>0\) et \(\|U_n-U_\star\|_2\) est suffisamment petit, alors \(U_n^\top U_n\) reste inversible et un calcul standard de perturbation matricielle donne
\[
\|P_{U_n}-P_{U_\star}\|_2
\le
C_1\|U_n-U_\star\|_2.
\]
Comme \(P_{U_n}=P_{T_n}\) et \(P_{U_\star}=P_0\), le résultat suit.
\end{proof}

\subsubsection{Théorème principal}

\begin{theorem}[Concentration du différentiel de projection]
Sous le modèle probabiliste ci-dessus, sous l'hypothèse de séparation géométrique, et pour \(n\) assez grand pour que la face active soit stable, on a :
\[
\mathbb P\!\left(
\|D\Pi_{M_n}(u^\star)-P_0\|_2
\le
C'\left(
\sqrt{\frac{k+\log(1/\delta)}{n}}
+
\sigma\sqrt{\frac{d+\log(1/\delta)}{n}}
\right)
\right)\ge 1-\delta,
\]
pour une constante \(C'>0\) dépendant uniquement de \(\tau\), \(\alpha\), et du conditionnement géométrique local de \(F_\star\).
\end{theorem}

\begin{proof}
Par le lemme de concentration uniforme du support, on a avec probabilité au moins \(1-\delta\),
\[
d_H(M_n,M_\star)
\le
C\left(
\sqrt{\frac{k+\log(1/\delta)}{n}}
+
\sigma\sqrt{\frac{d+\log(1/\delta)}{n}}
\right).
\]
Si \(n\) est assez grand pour que le membre de droite soit inférieur à \(c_0\), alors la face active reste stable. On peut donc écrire
\[
D\Pi_{M_n}(u^\star)=P_{T_n},
\qquad
D\Pi_{M_\star}(u^\star)=P_0.
\]
Le lemme de perturbation du projecteur tangent donne alors
\[
\|D\Pi_{M_n}(u^\star)-P_0\|_2
=
\|P_{T_n}-P_0\|_2
\le
C_1\, d_H(M_n,M_\star).
\]
En combinant les deux estimations, on obtient l'inégalité annoncée avec \(C'=C_1C\).
\end{proof}

\subsubsection{Corollaire : concentration du Hessien effectif}

\begin{corollary}[Hessien local]
Comme
\[
\nabla^2 d_{M_n}(u^\star)^2
=
2\bigl(I-D\Pi_{M_n}(u^\star)\bigr),
\]
on déduit
\[
\nabla^2 d_{M_n}(u^\star)^2
=
2(I-P_0)+O_{\mathbb P}\!\left(
\sqrt{\frac{k}{n}}+\sigma\sqrt{\frac{d}{n}}
\right).
\]
Par conséquent, pour l'énergie PCH,
\[
\nabla^2 J_{\mathrm{PCH}}(u^\star;w,c)
=
(1+2\lambda_{\mathrm{ridge}})I
+
2\lambda_{\mathrm{ctx}}(I-P_0)
+
O_{\mathbb P}\!\left(
\sqrt{\frac{k}{n}}+\sigma\sqrt{\frac{d}{n}}
\right).
\]
\end{corollary}

\subsubsection{Conséquence pour Cas 3}

Si, en plus, la famille de corpus, d'encodeurs, de réductions PCA et de réglages hyperparamétriques satisfait une hypothèse d'uniformité géométrique, au sens où les constantes de marge \(\tau,\alpha\) et le projecteur limite \(P_0\) sont uniformément contrôlés sur la famille considérée, alors
\[
D\Pi_{M_n}(u^\star)\xrightarrow{\mathbb P}P_0
\]
uniformément sur cette famille, et
\[
DT_\eta^{\mathrm{PCH}}(u^\star)
=
\bigl(I+\eta H_{\mathrm{eff}}\bigr)^{-1}
+
o_{\mathbb P}(1),
\]
avec
\[
H_{\mathrm{eff}}
=
(1+2\lambda_{\mathrm{ridge}})I
+
2\lambda_{\mathrm{ctx}}(I-P_0).
\]

\paragraph{Lecture.}
Le point crucial est bien le suivant : le différentiel de projection ne devient stable que si deux choses sont vraies simultanément :
\[
d_H(M_n,M_\star)\to 0
\quad\text{et}\quad
\text{la face active reste inchangée.}
\]
Sans la première, il n'existe aucune géométrie limite ; sans la seconde, le différentiel change de régime et aucune linéarisation stable n'est possible. Le Cas 3 est donc une conséquence de la concentration géométrique du hull et de la persistance des faces actives, et non un fait autonome du prox.



















































































%%%%%%%%%%%%%%%%%%%%%%%%%%%%%%%%%%%%%%%%%%%%%%%%%%%%%%%%%%%%%%%%%%%%%%%%%%%%%%%%%%%%%%%%%%%%%%%%%%%%%%%%%%%%%%%%%%%%%%%%%%%%%%%%%%%%%%%%%%%%%%%%%%%%%%%%%%%%%%%%%%%%%%%%%%%%%%%%%%%%%%%%%%%%%%%%%%%%%%%%%%%%%%%%%%%%%%%%%%%%%%%%%%%%%%%%%%%%%%%%%%%%%%%%%%%%%%%%%%%%%%%%%%%%%%%%%%%%%%%%%%%confrontation avec la theorie complete :

\section{Évaluation du niveau scientifique de la théorie complète et confrontation aux résultats expérimentaux}

\subsection{Préambule méthodologique}

L'évaluation présentée dans cette section porte sur la théorie complète décrite dans le document fondateur, et non uniquement sur le noyau PCH.

Une première lecture centrée exclusivement sur les expériences pourrait conduire à identifier PCH comme l'objet principal de la théorie. Une analyse intégrale du document montre cependant que le noyau PCH n'est qu'un composant d'une architecture théorique beaucoup plus vaste.

La théorie développée dans le document est une théorie \emph{PCH--ESM} complète comprenant plusieurs niveaux hiérarchiques :

\begin{enumerate}
    \item la théorie PH ;
    \item la théorie PCH ;
    \item la décomposition spectrale en espaces spécialisés ;
    \item le mécanisme de collage global \(G\) ;
    \item la théorie Multi-PH ;
    \item la modélisation hiérarchique de la croyance, de l'incertitude et de la représentation du monde.
\end{enumerate}

Par conséquent, la question scientifique correcte n'est pas :

\begin{quote}
``Les expériences valident-elles PCH ?''
\end{quote}

mais plutôt :

\begin{quote}
``Quelle partie de la théorie complète est démontrée analytiquement, quelle partie est soutenue expérimentalement, et quelle partie reste ouverte ?''
\end{quote}

\subsection{Valeur scientifique de la théorie seule}

On considère ici uniquement le contenu mathématique du document, indépendamment des expériences réalisées.

\subsubsection{Bloc PH}

La théorie PH constitue l'une des parties les plus solides de l'édifice.

Les résultats démontrés comprennent notamment :

\begin{itemize}
    \item existence du minimiseur ;
    \item unicité ;
    \item stabilité ;
    \item convergence proximale ;
    \item propriétés de contractivité ;
    \item existence des points fixes.
\end{itemize}

Ces résultats reposent essentiellement sur des outils standards d'analyse convexe, de théorie des opérateurs monotones et de théorie proximale.

Du point de vue mathématique, cette partie peut être considérée comme rigoureuse et mature.

\subsubsection{Bloc PCH}

Le noyau PCH hérite de la structure analytique de PH tout en introduisant une composante géométrique liée aux hulls convexes.

Le document établit notamment :

\begin{itemize}
    \item la purification contextuelle ;
    \item la définition variationnelle du noyau ;
    \item l'existence du minimiseur ;
    \item son unicité ;
    \item la stabilité de la dynamique associée ;
    \item les propriétés de convergence du schéma proximal.
\end{itemize}

Cette partie est également très solide du point de vue mathématique.

\subsubsection{Bloc de spécialisation spectrale}

Le document introduit une décomposition spectrale de la forme

\[
\mathcal H
=
\mathcal H_0
\oplus
\mathcal H_1
\oplus
\cdots
\oplus
\mathcal H_n.
\]

Cette construction repose sur des projecteurs spectraux et sur une décomposition orthogonale de l'espace latent.

Les résultats de cohérence mathématique apparaissent rigoureux.

Toutefois, contrairement à PH et PCH, cette partie dispose actuellement de peu de validations empiriques directes.

\subsubsection{Bloc de collage global \(G\)}

Le document introduit un mécanisme de reconstruction globale à partir des composantes spécialisées.

Les propriétés démontrées comprennent notamment :

\begin{itemize}
    \item cohérence du collage ;
    \item existence du minimiseur global ;
    \item unicité ;
    \item optimalité de type inverse-variance ;
    \item stabilité de la reconstruction.
\end{itemize}

Mathématiquement, cette partie apparaît cohérente et relativement développée.

\subsubsection{Bloc Multi-PH}

Le niveau Multi-PH constitue l'ambition la plus élevée du document.

Il introduit notamment :

\begin{itemize}
    \item des représentations hiérarchiques du monde ;
    \item des mécanismes de croyance ;
    \item des mécanismes d'incertitude ;
    \item des structures de composition entre espaces spécialisés ;
    \item une dynamique hiérarchique globale.
\end{itemize}

La formalisation mathématique est cohérente.

En revanche, cette partie ne bénéficie pas encore du même niveau de validation que PH ou PCH.

\subsection{Évaluation de la théorie seule}

Au vu des démonstrations disponibles, la hiérarchie suivante paraît raisonnable :

\begin{center}
\begin{tabular}{lc}
\hline
Bloc théorique & Niveau de maturité \\
\hline
PH & très élevé \\
PCH & très élevé \\
Spécialisation spectrale & élevé \\
Collage \(G\) & élevé \\
Multi-PH & moyen à élevé \\
Cas 3 universel & ouvert \\
\hline
\end{tabular}
\end{center}

\subsection{Apport des expériences réalisées}

Les campagnes expérimentales apportent une information nouvelle qui n'était pas contenue dans le document théorique seul.

Elles permettent une confrontation directe entre :

\begin{itemize}
    \item les prédictions analytiques ;
    \item les comportements numériques observés.
\end{itemize}

\subsection{Statut du Cas 1}

Les expériences montrent :

\begin{itemize}
    \item absence de violation énergétique ;
    \item convergence systématique ;
    \item résidus de point fixe compris entre \(10^{-17}\) et \(10^{-16}\).
\end{itemize}

Le Cas 1 est donc simultanément :

\[
\boxed{
\text{démontré analytiquement}
}
\]

et

\[
\boxed{
\text{validé expérimentalement}
}
\]

avec un niveau de précision exceptionnel.

\subsection{Statut du Cas 2}

La théorie prédit une réduction géométrique de la distance au hull admissible.

Les expériences montrent notamment :

\[
0.1532
\rightarrow
0.00746
\]

dans le cas oracle de l'expérience ex1.

La réduction est :

\begin{itemize}
    \item systématique ;
    \item forte ;
    \item stable ;
    \item reproductible.
\end{itemize}

Le Cas 2 doit donc être considéré comme

\[
\boxed{
\text{très fortement supporté expérimentalement}
}
\]

bien qu'il ne dispose pas encore d'une preuve universelle complète.

\subsection{Statut du Cas 3 local}

La découverte principale issue des expériences concerne la loi locale :

\[
\rho(\eta)
=
\frac{1}{1+\eta\mu_{\mathrm{eff}}}.
\]

Cette relation est observée avec des erreurs relatives de l'ordre de

\[
10^{-16}.
\]

Une telle précision est extrêmement rare dans un système réel impliquant :

\begin{itemize}
    \item des embeddings textuels ;
    \item des hulls convexes ;
    \item des solveurs numériques ;
    \item des données réelles.
\end{itemize}

Par conséquent :

\[
\boxed{
\text{Cas 3 local très fortement supporté}
}
\]

et probablement la découverte expérimentale la plus importante de l'ensemble du programme.

\subsection{Ce qui est réellement validé aujourd'hui}

Les expériences réalisées testent essentiellement :

\[
\boxed{
\text{PH}
\longrightarrow
\text{PCH}
}
\]

et très peu les autres composantes du système théorique.

En particulier, les expériences ne testent pratiquement pas :

\begin{itemize}
    \item la spécialisation spectrale complète ;
    \item la partition canonique ;
    \item le collage optimal \(G\) ;
    \item les compromis capacité/spécialisation ;
    \item Multi-PH ;
    \item la théorie hiérarchique de la représentation du monde.
\end{itemize}

Il serait donc scientifiquement incorrect d'affirmer aujourd'hui que la théorie complète est validée expérimentalement.

Ce qui est validé est essentiellement :

\[
\boxed{
\text{le bloc PH/PCH}
}
\]

et certaines conséquences locales associées.

\subsection{Évaluation globale}

\subsubsection{Théorie seule}

La théorie prise isolément présente :

\begin{itemize}
    \item une structure mathématique cohérente ;
    \item des définitions précises ;
    \item des théorèmes nombreux ;
    \item une architecture conceptuelle ambitieuse ;
    \item une absence de contradiction interne manifeste.
\end{itemize}

Son niveau scientifique peut être considéré comme élevé.

\subsubsection{Théorie plus expériences}

Les expériences apportent :

\begin{itemize}
    \item la validation du Cas 1 ;
    \item un fort soutien au Cas 2 ;
    \item une validation numérique extrêmement précise du Cas 3 local ;
    \item une caractérisation empirique de
    \[
    \mu_{\mathrm{eff}};
    \]
    \item l'identification du routage comme principal facteur limitant.
\end{itemize}

L'ensemble théorie + expériences possède donc une valeur scientifique significativement supérieure à celle du document seul.

\subsection{Résultat le plus important obtenu à ce jour}

Le résultat le plus important n'est probablement ni :

\begin{quote}
``PCH converge''
\end{quote}

ni

\begin{quote}
``PCH améliore l'accuracy''
\end{quote}

mais plutôt l'observation suivante :

\[
\boxed{
DT_\eta(u^\star)
\approx
(I+\eta H_{\mathrm{eff}})^{-1}
}
\]

avec une précision numérique proche de la précision machine.

Cette observation relie directement :

\begin{enumerate}
    \item la théorie variationnelle ;
    \item la géométrie des hulls convexes ;
    \item la dynamique locale du proximal ;
    \item les résultats observés sur des données réelles.
\end{enumerate}

Elle constitue actuellement le lien le plus fort entre la théorie abstraite et les observations expérimentales.

\subsection{Verdict final}

À l'heure actuelle, la situation scientifique peut être résumée ainsi :

\[
\boxed{
\text{PH : démontré et validé}
}
\]

\[
\boxed{
\text{PCH : démontré et validé}
}
\]

\[
\boxed{
\text{Cas 1 : confirmé}
}
\]

\[
\boxed{
\text{Cas 2 : fortement supporté}
}
\]

\[
\boxed{
\text{Cas 3 local : très fortement supporté}
}
\]

\[
\boxed{
\text{Cas 3 universel : ouvert}
}
\]

\[
\boxed{
\text{Spécialisation spectrale : peu testée}
}
\]

\[
\boxed{
\text{Collage }G\text{ : peu testé}
}
\]

\[
\boxed{
\text{Multi-PH : essentiellement non testé}
}
\]

La conclusion correcte n'est donc pas que la théorie complète est démontrée.

La conclusion correcte est que le bloc PH/PCH est désormais fortement étayé à la fois analytiquement et expérimentalement, tandis que les niveaux supérieurs de la théorie PCH--ESM restent encore largement à explorer et à confronter aux données.